% ch06-lambdatree.tex — λ-tree syntax: computing with binders (Miller & Nadathur, Chapter 7)
\chapter{\texorpdfstring{$\lambda$}{λ}-tree syntax: computing with binders}\label{ch:lambdatree}

The previous chapters gave us two kinds of machinery. \Cref{ch:lambda} and \cref{ch:hoprog} gave us
simply typed $\lambda$-terms as \emph{data}, with an equality that already contains $\alpha$-, $\beta$- and
$\eta$-conversion. \Cref{ch:hh} gave us the two goal connectives of hereditary Harrop formulas: the
universal goal \lp{pi x\ G}, proved with a fresh \emph{eigenvariable}, and the implicational goal
\lp{D => G}, proved with the program temporarily extended by \lp{D}. This chapter, which follows Chapter~7 of
the book \cite{MillerNadathur2012} (\book{175}--\book{210}), shows what the two kinds of machinery can do
together: they let a logic program compute over syntax that contains \emph{binders}.

This is the chapter that motivates the whole proposal of this document. Every program in it
manipulates formulas with quantifiers or $\lambda$-terms. None of them contains a single line about variable
names, renaming, or capture-avoiding substitution. All of that is \enquote{relegated to the meta-level}.
A miniKanren that wants to run these programs must therefore provide the meta-level: $\lambda$-terms as data,
eigenvariables with a scope discipline, and a scoped program extension. We will read the book's programs
closely, walk through how \lp{pi} and \lp{=>} make a binder \enquote{mobile}, and, in the \emph{Kanren
angle} boxes, record what each concept asks of \micro.

\section{The problem: syntax with binders}\label{sec:lambdatree-problem}

Many symbolic objects contain binders \book{175}--\book{176}: the $x$ in $\int f(x)\,dx$, in
$\frac{d}{dx}$ and in $\sum_{x=1}^{100}$; quantifiers; $\lambda$-abstractions; local declarations and formal
parameters in programming languages; type quantifiers of polymorphic types; name restriction in the
$\pi$-calculus. Programs that handle such objects (differentiators, compilers, interpreters, theorem
provers, type checkers) all need the same three notions \book{176}:
\begin{enumerate}[nosep]
  \item the distinction between \emph{free} and \emph{bound} occurrences of a variable;
  \item the irrelevance of bound names ($\alpha$-conversion): the formulas $\forall x.\,p(x)$ and
        $\forall y.\,p(y)$ are the same;
  \item \emph{capture-avoiding substitution}, needed both for instantiating quantifiers and for passing
        parameters.
\end{enumerate}
The book observes that \enquote{a computational treatment of such binding-related notions is complicated,
and programming language support for realizing them can be valuable} \book{176}. Anyone who has written a
substitution function over a named representation, with its renaming case, knows why.

The book's answer has three parts. We state them now and spend the chapter on them.
\begin{enumerate}[nosep]
  \item \textbf{Encoding.} Each object-level binder becomes a constant applied to a
        $\lambda$-abstraction of the meta-level. So \lp{all (x\ B)} encodes the formula $\forall x.\,B$, and
        \lp{abs (x\ t)} encodes the term $\lambda x.\,t$.
  \item \textbf{Substitution is $\beta$.} Instantiating the bound variable of an abstraction \lp{B} with
        \lp{T} is just writing the application \lp{(B T)}; the meta-level $\beta$-conversion does the rest.
  \item \textbf{Mobility.} To recurse \emph{under} a binder, the program turns the term-level $\lambda$ into a
        formula-level \lp{pi}, which proof search turns into a proof-level eigenvariable; facts about the
        new eigenvariable are assumed with \lp{=>}.
\end{enumerate}

\section{Encoding objects with binding structure}\label{sec:lambdatree-encoding}

\subsection{First-order formulas}\label{sec:lambdatree-fol}

The object logic is first-order logic over two sorts, \lp{term} and \lp{form}. The propositional
connectives are ordinary constants (re-listed on \book{177}):
\begin{lprolog}
type   ff,          % encoding the false proposition
       tt   form.   % encoding the true proposition
type   &&,          % encoding conjunction
       !!,          % encoding disjunction
       ==>  form -> form -> form.   % encoding implication
type   neg  form -> form.           % encoding negation
infixl &&  5.
infixl !!  4.
infixr ==> 3.
\end{lprolog}
A sample non-logical signature \book{177}:
\begin{lprolog}
type p  term -> form.
type q  term -> term -> form.
type f  term -> term.
type a  term.
\end{lprolog}
Terms of type \lp{term -> form} are \enquote{abstractions of a term over a formula} \book{178}, for example
\begin{lprolog}
x\ (p x) ==> (p (f x))
Z\ (p Z) ==> (p (f Z))
Z\ (p Z) && (p (f Z))
\end{lprolog}
The first two are \emph{the same} $\lambda$-term, since they differ only in the bound name. Note that the
capital \lp{Z} after \lp{Z\} is a bound variable, not a logic variable. The quantifiers are then
\emph{second-order constants}: their argument has a functional type \book{178}.
\begin{lprolog}
type all, some  (term -> form) -> form.
\end{lprolog}
\begin{lprolog}
all x\ (p x) ==> (p (f x))          % ∀x (p(x) ⊃ p(f(x)))
some Z\ (p Z) && (p (f Z))          % ∃Z (p(Z) ∧ p(f(Z)))
\end{lprolog}
The book states the principle behind this as follows: \enquote{the mechanism that we have described divides
the encoding of the object logic quantifiers into two parts: a second order constant that identifies the
quantifier and an abstraction that captures its binding effect} \book{178}. For instance
$\forall x\,\forall y\,(q(x,y)\imp\exists z\,(p(z)\wedge q(z,y)))$ is encoded as
\begin{lprolog}
all x\ all y\ (q x y) ==> some z\ p z && q z y.
\end{lprolog}
where \lp{\} extends as far to the right as possible and \lp{&&} (priority 5) binds tighter than
\lp{==>} (priority 3).

\begin{definition}[The encoding of formulas]\label{def:lambdatree-encoding}
For first-order terms and formulas over the signature above, define $\ulcorner\cdot\urcorner$ by
\[
\begin{aligned}
\ulcorner x\urcorner &= x, &
\ulcorner f(t_1,\dots,t_n)\urcorner &= f\,\ulcorner t_1\urcorner\cdots\ulcorner t_n\urcorner,\\
\ulcorner B\wedge C\urcorner &= \ulcorner B\urcorner \mathbin{\texttt{\&\&}} \ulcorner C\urcorner
\quad(\text{similarly for } \vee,\imp,\neg,\top,\bot), &
\ulcorner\forall x.\,B\urcorner &= \cst{all}\,(\abs{x}{\ulcorner B\urcorner}),\\
& & \ulcorner\exists x.\,B\urcorner &= \cst{some}\,(\abs{x}{\ulcorner B\urcorner}).
\end{aligned}
\]
\end{definition}

The book does not state $\ulcorner\cdot\urcorner$ as a definition or prove anything about it; this formal
version is ours, and so is the name. The chapter relies on two properties of it without stating them.
First, it is a bijection between closed object formulas over the signature, taken up to renaming of bound
variables, and closed terms of type \lp{form} over the same signature, taken up to $\alpha\beta\eta$-conversion.
A canonical representative of each class is the $\beta$-normal, $\eta$-long term, which is what
$\ulcorner\cdot\urcorner$ produces. Second, it commutes with substitution:
$\ulcorner \subst{B}{t}{x}\urcorner = \subst{\ulcorner B\urcorner}{\ulcorner t\urcorner}{x}$. The second
property is what makes \enquote{substitution is $\beta$} correct.

\subsection{Untyped \texorpdfstring{$\lambda$}{λ}-terms}\label{sec:lambdatree-untyped}

The second running object language is the untyped $\lambda$-calculus \book{178}--\book{179}:
\begin{lprolog}
kind tm   type.
type app  tm -> tm -> tm.
type abs  (tm -> tm) -> tm.
\end{lprolog}
\begin{center}
\begin{tabular}{@{}ll@{}}
\toprule
untyped term & encoding \\
\midrule
$\lambda x\, x$ & \lp{(abs x\x)} \\
$\lambda x\,(x\,x)$ & \lp{(abs x\ (app x x))} \\
$\lambda x\,\lambda y\, x$ & \lp{(abs x\ abs y\ x)} \\
$\lambda x\,\lambda y\, y$ & \lp{(abs x\ abs y\ y)} \\
$\lambda x\,\lambda y\,(y\,x)$ & \lp{(abs x\ abs y\ app y x)} \\
$\lambda x\,\lambda y\,\lambda z\,((x\,z)\,(y\,z))$ & \lp{(abs x\abs y\abs z\ (app (app x z) (app y z)))} \\
$(\lambda x\,(x\,x))\,(\lambda x\,(x\,x))$ & \lp{(app (abs x\ app x x) (abs x\ app x x))} \\
\bottomrule
\end{tabular}
\end{center}
Each object application becomes one \lp{app}; each object $\lambda$ becomes one \lp{abs} followed by one
meta-level $\lambda$. Two remarks help in reading the rest of the chapter.
\begin{itemize}[nosep]
  \item No closed term of type \lp{tm} can be built from \lp{app} alone: every closed term starts with an
        \lp{abs}, whose argument (of type \lp{tm -> tm}) is something like \lp{x\x} or \lp{x\ app x x}.
  \item \enquote{the type tm can be seen as admitting new constructors as one descends under binders}
        \book{179}. Below \lp{abs x\ abs y\} both \lp{x} and \lp{y} behave as extra constants of type
        \lp{tm}. Mobility (\cref{sec:lambdatree-mobility}) makes this literally true: under a binder the
        program works with a genuinely new constant.
\end{itemize}

\subsection{What the encoding gives for free}\label{sec:lambdatree-properties}

The abstractions over syntax live at arrow types (\lp{term -> form}, \lp{tm -> tm}), and the equality of
those types includes the three conversions of \cref{ch:lambda}.
\begin{description}[nosep,font=\normalfont\itshape]
  \item[$\alpha$.] The object language identifies formulas up to renaming of bound variables, and so does
        the meta-language. \enquote{Since $\alpha$-conversion is included within the logical notion of equality,
        it is impossible to access the names of bound variables} \book{180}. The book presents this as a
        strength: the programmer is spared what does not matter, and the implementer may drop names entirely
        (for instance with de Bruijn indices \cite{deBruijn1972}). The cost is that printing and parsing must
        invent names.
  \item[$\beta$.] This is the subject of \cref{sec:lambdatree-beta}.
  \item[$\eta$.] Its effect is \enquote{rather mild} \book{180}: for $f:\typ{i}\to\typ{i}$,
        $f=_\eta\abs{x}{f\,x}$. Hence $\forall x\,(p\,x)$ may be written \lp{(all x\ p x)} \emph{or}
        \lp{(all p)}: \enquote{not only does the hohh language abstract away from names of bindings, but it also
        sometimes allows one to abstract away from the presence of actual binders} \book{180}.
\end{description}

\begin{pitfall}[An $\eta$-short quantifier is still a quantifier]
A head pattern \lp{(all B)} must also match the term \lp{(all p)}, and \lp{(B T)} must then reduce to
\lp{(p T)}. An implementation either unifies modulo $\eta$ or keeps terms $\eta$-long. A representation that
stores the binder \enquote{shape} syntactically and expects every quantifier to carry an explicit
$\lambda$ will reject \lp{(all p)}.
\end{pitfall}

\section{Object-level substitution is meta-level \texorpdfstring{$\beta$}{β}}\label{sec:lambdatree-beta}

\subsection{Instantiating quantifiers}

The formula $\forall x\,(p\,x\imp\exists y\,(q\,x\,y))$ is \lp{all x\ p x ==> some y\ q x y}. To instantiate
it with $f(a)$ we simply apply the abstraction to \lp{(f a)} \book{180}--\book{181}:
\begin{lprolog}
(x\ p x ==> some y\ q x y) (f a)
\end{lprolog}
which $\beta$-converts to \lp{p (f a) ==> some y\ q (f a) y}. (The book prints \lp{p (h a)} here, a typo;
the page image confirms it, and it must be \lp{p (f a)}.) Note what was \emph{not} needed: no renaming of
the inner \lp{y}, no check that \lp{(f a)} has no free \lp{y}. Capture avoidance is the meta-level's job.

The predicate \lp{list_instan} relates a list of $n\ge 0$ terms and a formula with at least $n$ outer
universal quantifiers to the result of instantiating the first $n$ \book{181}:
\begin{lprolog}
type list_instan  list term -> form -> form -> o.
list_instan nil B B.
list_instan (T::Ts) (all B) C :- list_instan Ts (B T) C.
\end{lprolog}
The head \lp{(all B)} is $\eta$-equal to \lp{(all x\ B x)}: matching it binds \lp{B} to the abstraction
itself, and \enquote{the substitution for a bound variable is realized simply by writing down (B T)}
\book{181}.

\subsection{An interpreter for an object logic}\label{sec:lambdatree-interp}

A more substantial example is an interpreter for first-order Horn clauses, written in higher-order Horn
clauses (Figure~7.1, \book{181}). The object-level program is an explicit argument \lp{D}, a single formula.
\begin{lprolog}
type interp     form -> form -> o.
type backchain  form -> form -> form -> o.
type atom       form -> o.

interp D tt.
interp D (G1 && G2)  :- interp D G1, interp D G2.
interp D (G1 !! G2)  :- interp D G1; interp D G2.
interp D (some   G)  :- interp D (G X).
interp D A           :- atom A, backchain D D A.
backchain D A A.
backchain D (D1 && D2) A   :- backchain D D1 A; backchain D D2 A.
backchain D (all D1) A     :- backchain D (D1 X) A.
backchain D (G ==> D1) A   :- backchain D D1 A, interp D G.
\end{lprolog}
The two higher-order clauses are the interesting ones. An object goal \lp{(some G)} is solved by solving
\lp{(G X)} for a \emph{fresh logic variable} \lp{X}: the meta-level $\beta$ plugs the logic variable in for the
object-bound variable, and later unification will find a witness. An object clause \lp{(all D1)} is used by
instantiating it with a fresh logic variable as well. The sample object program of Figure~7.2 (\book{182})
is
\begin{lprolog}
type a, b, c    term.
type adj, path  term -> term -> form.
type prog       form -> o.
atom (adj _ _) & atom (path _ _).
prog ((adj a b) && (adj b c) &&
      (all x\ all y\ (adj x y) ==> (path x y)) &&
      (all x\ all y\ all z\ (adj x y) && (path y z) ==> (path x z))).
\end{lprolog}
Two notational points recur in the rest of the chapter. An \lp{&} between two clause heads, as in
\lp{atom (adj _ _) & atom (path _ _).}, is clause-level conjunction. The heads share one body (here an empty
one), so the line stands for two clauses written together. The underscore \lp{_} is an anonymous logic
variable. The predicate \lp{atom} is not part of the interpreter. It is defined for each object signature by
listing its predicate constants, and here it says that \lp{adj} and \lp{path} atoms are atomic.
The query \lp{?- prog P, interp P (path a X).} has two answers, \lp{X = b} and then \lp{X = c} \book{182}.

\begin{tutorial}[How \texttt{interp} finds \texttt{X = b}]
The goal \lp{interp P (path a X)} matches the last \lp{interp} clause, because \lp{atom (path a X)} holds. We
must prove \lp{backchain P P (path a X)}. The program \lp{P} is a left-nested \lp{&&} of four parts, so the
second \lp{backchain} clause descends into it. The two \lp{adj} facts do not match \lp{path a X}. The third part
is \lp{all x\ all y\ (adj x y) ==> (path x y)}: the \lp{(all D1)} clause applies it to a fresh \lp{X1}, giving
\lp{(D1 X1)} $=_\beta$ \lp{all y\ (adj X1 y) ==> (path X1 y)}; once more with \lp{Y1}, giving
\lp{(adj X1 Y1) ==> (path X1 Y1)}. The \lp{==>} clause then requires \lp{backchain P (path X1 Y1) (path a X)}
(the fact clause, binding \lp{X1 = a}, \lp{Y1 = X}) and then \lp{interp P (adj a X)}, which backchains to the fact
\lp{adj a b} with \lp{X = b}. On backtracking, the fourth part of the program yields \lp{X = c} through
\lp{adj a b} and \lp{path b c}. At no point did the interpreter rename a bound variable: each $\beta$-step
\lp{(D1 X1)} replaced an object-level bound variable by a logic variable.
\end{tutorial}

\subsection{Evaluators for the untyped \texorpdfstring{$\lambda$}{λ}-calculus}

Call-by-name and call-by-value evaluation are four clauses (Figure~7.3, \book{182}):
\begin{lprolog}
type cbn, cbv  tm -> tm -> o.
cbn (abs R) (abs R).
cbn (app M N) V :- cbn M (abs R), cbn (R N) V.
cbv (abs R) (abs R).
cbv (app M N) V :- cbv M (abs R), cbv N U, cbv (R U) V.
\end{lprolog}
\lp{(R N)} is the object-level $\beta$-contraction, done by the meta-level $\beta$. (The text on \book{182}
says \enquote{the term U such that eval N U is provable}; \lp{eval} is a loose name for \lp{cbv}.) The behaviour
is the expected one \book{183}. $\Omega=(\lambda x\,(x\,x))(\lambda x\,(x\,x))$ has no value under either
evaluator: the search for a derivation is infinite. On $(\lambda x\,\lambda w\,w)\,\Omega$, call-by-name
terminates and call-by-value loops:
\begin{lprolog}
?- cbn (app (abs x\ abs w\w) (app (abs x\ app x x) (abs x\ app x x))) V.
   % V = (abs w\w)
?- cbv (app (abs x\ abs w\w) (app (abs x\ app x x) (abs x\ app x x))) V.
   % infinite search
\end{lprolog}
The book draws the moral of this section in one paragraph \book{183}. The specifications are succinct and
declarative because \enquote{all the details regarding object-level bound variables and object-level
substitution has been relegated to the meta-level}. They are still formal and precise: the details
\enquote{have not gone away; they simply have become the burden of an implementation of a metalogic such as
$\lambda$Prolog}, and the programmer \enquote{is relieved of dealing with the details behind such
\textquote{concrete nonsense} as binder names, capture avoiding substitutions, etc.}

\begin{kanren}[What works today: abstractions as Scheme closures]
All the programs of this section use an abstraction only as an \emph{input}. The abstraction is already known
when \lp{(R N)} is formed. A miniKanren can imitate this with \emph{functional} higher-order abstract syntax:
represent \lp{abs R} by a list holding a Scheme procedure, and perform \lp{(R N)} with \scm{project°}. This
runs on today's \micro, unchanged:
\begin{scheme}
(define (abs f) `(abs ,f))
(define (app m n) `(app ,m ,n))
(define ω (abs (λ (x) (app x x))))
(define Ω (app ω ω))
(define K* (abs (λ (x) (abs (λ (w) w)))))          ; λx λw w

; call-by-name evaluation, Figure 7.3: (R N) is Scheme application, available only through project°
(define-relation (cbn° m v)
  (cond°
    ((fresh° (r) (=° m `(abs ,r)) (=° v m)))
    ((fresh° (m1 n r) (=° m `(app ,m1 ,n)) (cbn° m1 `(abs ,r))
       (project° ((r* r)) (cbn° (r* n) v))))))

; printing a HOAS term needs invented names: apply each closure to a generated symbol
(define (hoas->sexp t)
  (let loop ((t t) (k 1))
    (cond
      ((and (pair? t) (eq? (car t) 'abs))
       (let ((w (string->symbol (string-append "w" (number->string k)))))
         `(λ (,w) ,(loop ((cadr t) w) (+ k 1)))))
      ((and (pair? t) (eq? (car t) 'app)) `(,(loop (cadr t) k) ,(loop (caddr t) k)))
      (else t))))

(print (μkanren-run (q -1 #t) (fresh° (v) (cbn° (app K* Ω) v) (project° ((v* v)) (=° q (hoas->sexp v*))))))
; α-equivalent closures are different Scheme objects: unification fails
(print (μkanren-run (q -1 #t) (=° (abs (λ (x) x)) (abs (λ (y) y)))))
\end{scheme}
\begin{answers}
((λ (w1) w1))
()
\end{answers}
The first answer is the book's \lp{V = (abs w\w)}. The second answer is the problem. Two $\alpha$-equivalent
abstractions are distinct procedures, and \micro{} compares procedures with \scm{eq?}, so they do not unify. A
closure cannot be printed without being applied to invented names (that is what \scm{hoas->sexp} does). A
closure cannot be inspected except by applying it, and no unification can \emph{produce} one. Recognisers
such as \lp{term}, \lp{bnorm} and \lp{typeof} below could still be run on closures. They only apply an
input abstraction to a fresh constant, although they also need \lp{pi} and \lp{=>}. But every program that
\emph{builds} an abstraction as output (\lp{red1}, path synthesis, \lp{trans} backwards, \lp{nnf},
\lp{prenex}, \lp{copy}, \lp{subst} backwards), and every program that asks whether a body depends on its
bound variable (\lp{vacuous}, the $\eta$-redex), is out of reach. \Cref{sec:lambdatree-hoas} comes back to
this \enquote{functional} reading of higher-order abstract syntax.
\end{kanren}

\section{Mobility of binders}\label{sec:lambdatree-mobility}

\subsection{Why an abstraction cannot be taken apart}

Equal expressions should have equal parts: equal non-empty lists have equal heads and equal tails. Applied to
$\lambda$-terms, this principle says that \enquote{we should not be able to decompose an abstraction into the
variable that it binds and its body} \book{183}. The terms \lp{(all x\ some y\ p x ==> p y)} and
\lp{(all z\ some y\ p z ==> p y)} are equal, yet the names \lp{x} and \lp{z} differ, and so do the would-be
bodies \lp{(some y\ p x ==> p y)} and \lp{(some y\ p z ==> p y)}. \enquote{Such decompositions, in fact, have no
logical status within the hohh language: Because equality includes $\alpha$-conversion, the name for a bound
variable is a fiction, and an implementation may not represent it at all} \book{183}.

So how does a program recurse into the body of a binder? The book's answer \book{184}:
\enquote{during proof search, a term-level binding (i.e., a $\lambda$-abstraction) can be converted into a
formula-level binding (i.e., a quantifier), which then can be converted into a proof-level binding (i.e., an
eigenvariable).}
\[
\underbrace{\abs{y}{t}}_{\text{term level}}
\;\rightsquigarrow\;
\underbrace{\forall x.\,(\cdots\imp P\,((\abs{y}{t})\,x))}_{\text{formula level}}
\;=_\beta\;
\forall x.\,(\cdots\imp P\,(\subst{t}{x}{y}))
\;\rightsquigarrow\;
\underbrace{P\,(\subst{t}{c}{y})\ \text{ with } c \text{ fresh}}_{\text{proof level}}
\]

\subsection{The proof-search rules behind mobility}

Recall from \cref{ch:hh} the rules for the two connectives that make mobility work, together with the
backchaining rule. A goal is proved relative to a signature $\Sig$ (constants and eigenvariables in scope)
and a program $\Prog$.
\begin{mathpar}
\inferrule{\seq{\Sig\cup\{c\}}{\Prog}{\subst{G}{c}{x}} \\ c\notin\Sig}{\seq{\Sig}{\Prog}{\forall x.\,G}}\;\rulename{generic}
\and
\inferrule{\seq{\Sig}{\Prog\cup\{D\}}{G}}{\seq{\Sig}{\Prog}{D\imp G}}\;\rulename{augment}
\and
\inferrule{\seq{\Sig}{\Prog}{\subst{G}{t}{x}} \\ t \text{ a } \Sig\text{-term}}{\seq{\Sig}{\Prog}{\exists x.\,G}}\;\rulename{instance}
\end{mathpar}
For an atomic goal $A$, \rulename{backchain} picks a clause $\forall\bar x\,(G\imp A')$ of $\Prog$ and an instance
with $A'\theta\eqb A$, then proves $G\theta$. In an implementation the term $t$ of \rulename{instance} and the
substitution $\theta$ are found by unification with logic variables. Then the side condition $c\notin\Sig$
becomes the \textbf{scoping constraint}: a logic variable created \emph{before} the eigenvariable $c$ must never
be instantiated with a term containing $c$, because at the time it was created its witness had to be a
$\Sig$-term. This is why \lp{sigma X\ pi y\ X = y} fails while \lp{pi y\ sigma X\ X = y} succeeds. Chapter~7
never mentions the constraint, but several of its programs (\lp{vacuous}, the $\eta$-redex clause,
\lp{prenex}) are correct only because of it.

\subsection{A first example: recognising terms}\label{sec:lambdatree-term}

Because of typing, the predicate \lp{term : tm -> o} could be defined by the single clause \lp{term T.}. A more
instructive definition reflects the two constructors of \lp{tm} into two clauses \book{184}. Written so that
their logical content is visible:
\begin{lprolog}
pi M\ term M => pi N\ term N => term (app M N).
pi R\ (pi x\ term x => term (R x)) => term (abs R).
\end{lprolog}
or, equivalently, in the usual clause syntax:
\begin{lprolog}
term (app M N) :- term M, term N.
term (abs R)   :- pi x\ term x => term (R x).
\end{lprolog}

\begin{tutorial}[Following the binder of \texttt{(abs y\textbackslash\ app y y)}]
The book traces the query \lp{?- (term (abs y\ app y y)).} \book{184}--\book{185}.
\begin{lprolog}
?- (term (abs y\ app y y)).
?- pi x\ term x => term ((y\ app y y) x).     % backchain on abs clause
?- pi x\ term x => term (app x x).            % β-conversion: term-level y -> formula-level x
?- term (app c c).                            % eigenvariable c, program + (term c)
\end{lprolog}
\begin{enumerate}[nosep]
  \item \emph{Backchaining} on the \lp{abs} clause unifies \lp{R} with the abstraction \lp{y\ app y y}. Nothing
        has been decomposed: \lp{R} holds the whole abstraction.
  \item The body of the clause applies \lp{R} to the \lp{pi}-bound \lp{x}. The redex \lp{((y\ app y y) x)} is
        a $\beta_0$-redex (\cref{sec:lambdatree-llambda}): a $\lambda$ applied to a bound variable. Contracting
        it moves the term-level binder \lp{y} to the formula-level binder \lp{x}.
  \item \rulename{generic} replaces \lp{x} by a new eigenvariable \lp{c}. \rulename{augment} adds the clause
        \lp{term c} to the program. The formula-level binder has become a proof-level binder.
  \item \lp{term (app c c)} backchains on the \lp{app} clause, and both subgoals \lp{term c} are closed by the
        assumed clause.
\end{enumerate}
Here \lp{c} is exactly one of the \enquote{new constructors} of \lp{tm} that appear under a binder. The program
has no clause for variables, and it needs none: each variable gets its clause at the moment its binder is
crossed, and loses it when the proof of the body is complete.
\end{tutorial}

The book distils the example into a programming idiom. We record it as a definition because every program
of the chapter is an instance of it.

\begin{definition}[The binder-mobility idiom, \book{185}]\label{def:lambdatree-idiom}
\enquote{In order to continue a recursive analysis within the scope of a binder, first, apply that binder to a
universal quantified goal variable, for which the metalogic will substitute an eigenvariable, and second,
using an implicational goal, assume new clauses that extend the definition of various predicates so that
they can deal with the presence of this new eigenvariable.}
\end{definition}

In code, the idiom is always the shape
\lp{p (c R) ... :- pi x\ (hyp x => p (R x) ...).} The abstraction \lp{R} is never opened. It is applied to a
fresh eigenvariable, and the assumed \lp{hyp x} tells the predicates of the program what to do when they meet
it.

\begin{pitfall}[A fresh constant is not an eigenvariable]
Proving \lp{pi x\ G} with a fresh constant is only half of \rulename{generic}. The other half is the scoping
constraint. Here is what happens without it, in today's \micro, if a universal goal is \enquote{implemented}
by handing a \scm{gensym} to the body:
\begin{scheme}
; a NAIVE universal goal: a fresh Scheme constant, with no eigenvariable scope check
(define (naive-pi f) (f (gensym 'c)))

; λProlog: ?- sigma X\ pi y\ X = y.   must fail; the naive encoding succeeds
(print (μkanren-run (q -1 #t) (naive-pi (λ (y) (=° q y)))))
; λProlog: ?- pi y\ sigma X\ X = y.   must succeed (X is created inside the scope of y)
(print (μkanren-run (q -1 #t) (naive-pi (λ (y) (fresh° (x) (=° x y) (=° q 'ok))))))
\end{scheme}
\begin{answers}
(c102)
(ok)
\end{answers}
The first query should fail. The query variable \scm{q} was created before the constant, yet it now holds that
constant, which has escaped its scope. The same mistake in a unifier, letting a clause variable be bound to a
term that mentions a variable bound inside its scope, makes \lp{vacuous} (\cref{sec:lambdatree-fol-comp})
accept formulas whose quantifier is not vacuous. The constraint is not a matter of efficiency: it is part of
the meaning of $\forall$.
\end{pitfall}

\begin{kanren}[What mobility requires of a miniKanren]
Mobility needs three things that \micro{} lacks (see \cref{ch:micro} for the engine).
\begin{enumerate}[nosep]
  \item \emph{Eigenvariables}: constants created during the search, distinct from every other term, with a
        scope. The engine must enforce the scoping constraint: a logic variable older than an eigenvariable
        can never be bound to a term containing it, \emph{even indirectly} through younger variables (binding
        $X:=f(Z)$ and later $Z:=c$ leaks $c$ into $X$). λProlog implementations use universe levels
        \cite{NadathurJayaramanKwon1995}. The miniKanren with symbolic constraints uses an absento-like
        \enquote{eigen} constraint over the variables in scope \cite{HemannByrdFriedman2013}.
  \item \emph{Scoped program extension}: \lp{D => G} adds \lp{D} only while \lp{G} is being proved. Other branches
        of the search and the answers of the enclosing goal must not see \lp{D}. Because miniKanren's search
        interleaves branches, the extension must live in the search \emph{state}, not in a host-language
        dynamic variable. A relation must be able to consult the assumed clauses as well as its static ones.
  \item \emph{Application of an abstraction to an eigenvariable}, \lp{(R x)}, on both sides of the equation:
        as input (when \lp{R} is known) and as output (when \lp{R} is a logic variable to be solved, as in most
        programs below).
\end{enumerate}
How the proposal provides each of these is the subject of \cref{ch:design}.
\end{kanren}

\section{Computing with untyped \texorpdfstring{$\lambda$}{λ}-terms}\label{sec:lambdatree-untyped-comp}

This section applies the idiom to the untyped $\lambda$-calculus. The programs grow from pure recognisers
to programs that build abstractions as output, to programs that assume rules with bodies.

\subsection{\texorpdfstring{$\beta$}{β}-normal forms}\label{sec:lambdatree-bnorm}

\begin{definition}[$\beta$-normal, $\beta$-body-normal, \book{185}]\label{def:lambdatree-bnorm}
An untyped $\lambda$-term is in \emph{$\beta$-normal form} if it contains no $\beta$-redex. Equivalently and
positively, it has the shape $\lambda x_1\ldots\lambda x_n\,(v\,t_1\ldots t_m)$ with $n,m\ge0$, where $v$ is a
variable and every $t_i$ is $\beta$-normal. The inner part $v\,t_1\ldots t_m$ is called
\emph{$\beta$-body-normal}.
\end{definition}

The definition becomes three inference rules and three clauses \book{186}:
\begin{mathpar}
\inferrule{\forall x.\ \mathrm{bbnorm}(x)\imp\mathrm{bnorm}(M\,x)}{\mathrm{bnorm}(\cst{abs}\,M)}
\and
\inferrule{\mathrm{bbnorm}(H)}{\mathrm{bnorm}(H)}
\and
\inferrule{\mathrm{bbnorm}(M) \\ \mathrm{bnorm}(N)}{\mathrm{bbnorm}(\cst{app}\,M\,N)}
\end{mathpar}
\begin{lprolog}
type bnorm, bbnorm  tm -> o.

bnorm (abs M)    :- pi x\ bbnorm x => bnorm (M x).
bnorm H          :- bbnorm H.
bbnorm (app M N) :- bbnorm M, bnorm N.
\end{lprolog}
There is no static clause saying that a variable is $\beta$-body-normal, and there could not be one, since
bound variables are not terms of the program. A variable becomes body-normal only through the assumption
\lp{bbnorm x} made when its binder is crossed. A redex \lp{app (abs R) N} is rejected because
\lp{bbnorm (abs R)} has no clause.

\begin{tutorial}[A normal term and a redex]
This walk-through is ours; the book gives no query for \lp{bnorm}. Take first
\lp{?- bnorm (abs x\ app x x).}
\begin{enumerate}[nosep]
  \item The \lp{abs} clause matches with \lp{M = x\ app x x}. The body is
        \lp{pi x\ bbnorm x => bnorm (M x)}: an eigenvariable \lp{c} is created, \lp{bbnorm c} is assumed, and the
        goal is \lp{bnorm (app c c)} (after $\beta_0$).
  \item \lp{bnorm (app c c)} is not an \lp{abs}, so the second clause reduces it to \lp{bbnorm (app c c)}, and the
        third clause to \lp{bbnorm c} and \lp{bnorm c}.
  \item \lp{bbnorm c} is the assumption. \lp{bnorm c} goes through the second clause to \lp{bbnorm c} again, also
        the assumption. The query succeeds.
\end{enumerate}
Now take \lp{?- bnorm (app (abs x\x) (abs x\x)).} The \lp{abs} clause does not match. The second clause asks for
\lp{bbnorm (app (abs x\x) (abs x\x))}, and the third clause for \lp{bbnorm (abs x\x)}. No static clause has a
head \lp{bbnorm (abs _)}, and the only assumptions ever made are \lp{bbnorm c} on eigenvariables, so the query
fails. The structure of the clauses alone, \enquote{a variable is body-normal only when its binder said so},
rules the redex out.
\end{tutorial}

A pleasant consequence is a substitution lemma that costs nothing.

\begin{theorem}[Substitution lemma \enquote{for free}, \book{186}]\label{thm:lambdatree-bnorm-subst}
Let $\Gamma=\{\cst{bbnorm}\,c_1,\ldots,\cst{bbnorm}\,c_n\}$, $n\ge0$, with the $c_i$ eigenvariables. If
\lp{(bnorm (abs R))} and \lp{(bbnorm S)} are both derivable from the program above together with $\Gamma$, then
\lp{(bnorm (R S))} is derivable too. In object-level terms: if $\lambda x\,t$ is $\beta$-normal and $s$ is
$\beta$-body-normal, then $\subst{t}{s}{x}$ is $\beta$-normal.
\end{theorem}
\begin{proof}[Proof (book's argument, with the inversion step made explicit)]
A derivation of \lp{bnorm (abs R)} must end with the \lp{abs} clause, so
\lp{pi x\ bbnorm x => bnorm (R x)} is derivable. Derivability is intuitionistic, so the universal quantifier
can be instantiated with \lp{S}, giving \lp{bbnorm S => bnorm (R S)}, and modus ponens with \lp{bbnorm S}
gives \lp{bnorm (R S)}.
\end{proof}
The book says only \enquote{since derivability is based on intuitionistic logic}. Two facts sit behind that
phrase: cut admissibility, and the fact that an eigenvariable may be instantiated by any term. The inversion
step (\enquote{must end with the \lp{abs} clause}) also uses the shape of $\Gamma$. $\Gamma$ contains only
atoms \lp{bbnorm c} on eigenvariables, so no assumption can match \lp{bnorm (abs R)}. In a named
representation this lemma is a page of index bookkeeping. Here the logic proves it.

\subsection{Reduction to normal form}\label{sec:lambdatree-reduce}

Figure~7.4 (\book{187}) specifies one-step reduction and reduction to normal form:
\begin{lprolog}
type redex, red1, reduce  tm -> tm -> o.
redex (app (abs R) N) (R N).
red1 M N :- redex M N.
red1 (app M N) (app P N) & red1 (app N M) (app N P) :- red1 M P.
red1 (abs M) (abs N) :- pi x\ red1 (M x) (N x).
reduce M M :- bnorm M.
reduce M N :- red1 M P, reduce P N.

type repeat   (A -> A -> o) -> A -> A -> o.
type reduce'  tm -> tm -> o.
repeat Pred M N :- Pred M P, !, repeat Pred P N.
repeat Pred M M.
reduce' M N :- repeat red1 M N.
\end{lprolog}
\begin{itemize}[nosep]
  \item \lp{redex} relates a $\beta$-redex to its contractum, again by meta-level $\beta$: \lp{(R N)}.
  \item \lp{red1} contracts exactly one redex somewhere in the term. The \lp{abs} clause is new. Its
        \emph{output} is the abstraction \lp{N}, and the goal \lp{red1 (M x) (N x)} asks unification to solve
        \lp{(N x) = t} for an unknown \lp{N}: the answer is \lp{N} $:=$ \lp{x\ t}, with the eigenvariable abstracted
        out of \lp{t}. This is a \emph{higher-order pattern} problem (\cref{sec:lambdatree-llambda,ch:hounif}).
        There is no \lp{=>} here: \lp{red1} should simply fail on a bare variable, so no clause about \lp{x} is
        needed.
  \item \lp{reduce} uses \lp{bnorm} to decide when to stop. \lp{reduce'} uses the polymorphic higher-order
        \lp{repeat}, which iterates its argument until it fails. The cut \lp{!} \enquote{prune[s] away alternative
        solutions to the parent goal} \book{186}. (The text there says \lp{repeat'}; the figure's predicate is
        \lp{repeat}.)
\end{itemize}

\begin{tutorial}[One step under a binder: \texttt{red1 (abs x\textbackslash\ app (abs y\textbackslash y) x) N}]
This is our own example of an output abstraction, traced against the clauses of Figure~7.4 without the
$\eta$-clause below. The term is $\lambda x\,((\lambda y\,y)\,x)$, with a redex under the binder.
\begin{enumerate}[nosep]
  \item \lp{red1 M N :- redex M N} fails: the term is an \lp{abs}, not an \lp{app}. The two \lp{app} clauses do
        not match either.
  \item The \lp{abs} clause matches with \lp{M = x\ app (abs y\y) x} and \lp{N = abs N1} for a new logic
        variable \lp{N1 : tm -> tm}. It creates an eigenvariable \lp{c}. The goal is
        \lp{red1 (app (abs y\y) c) (N1 c)}.
  \item The \lp{redex} clause applies with \lp{R = y\y} and contractum \lp{(R c)}, which is \lp{c}. Now unification
        must solve \lp{(N1 c) = c}. This is a pattern equation: \lp{N1} is applied to the eigenvariable \lp{c}
        only. Its unique solution abstracts \lp{c} out of the right-hand side: \lp{N1 := x\ x}.
  \item The answer is \lp{N = abs x\x}. On backtracking, the \lp{app} clauses try to reduce \lp{(abs y\y)} or
        \lp{c} in one step. The first descends under \lp{y\} to the goal \lp{red1 d _} on a new eigenvariable
        \lp{d}, the second gives \lp{red1 c _}. Both fail, because no clause, static or assumed, applies to a
        bare eigenvariable. There is no other answer.
\end{enumerate}
The program contains no code for \enquote{rebuild the binder around the reduced body}. Unification does it in
step 3, and the scoping constraint guarantees that \lp{N1}, created before \lp{c}, receives \lp{c} only through
its argument.
\end{tutorial}

The $\eta$-redex needs one clause \book{187}:
\begin{lprolog}
redex (abs x\ app M x) M.
\end{lprolog}
It is correct because of the scoping constraint. \lp{M} is a clause variable, quantified \emph{outside} the
binder \lp{x}, so no unifier can let \lp{M} depend on \lp{x}. This is exactly the side condition of $\eta$:
$x$ not free in $M$. For instance, \lp{abs x\ app x x} is not an $\eta$-redex: matching it against
\lp{abs x\ app M x} would need \lp{M = x}, and the scoping constraint forbids that binding. With this clause
added, the term of the tutorial above is also an $\eta$-redex. Matching \lp{abs x\ app (abs y\y) x} against
\lp{abs x\ app M x} gives \lp{M = abs y\y}, which does not mention \lp{x}, so \lp{red1} now finds
\lp{abs y\y} first, through \lp{redex}, and then the $\alpha$-equal \lp{abs x\x} again through the \lp{abs}
clause.

The book adds a semantic remark \book{186}--\book{187}. A model of the untyped $\lambda$-calculus is a domain
$D$ with maps $f:D\to[D\to D]$ and $g:[D\to D]\to D$; \lp{tm} plays $D$, \lp{app} plays $f$, \lp{abs} plays
$g$. The $\beta$-clause states $(f\circ g)\,R=R$ and the $\eta$-clause states $(g\circ f)\,M=M$: \enquote{these
are the familiar properties stating that f and g are, in fact, retracts}.

\subsection{Reduction based on paths}\label{sec:lambdatree-paths}

A \emph{path} walks from the root of a term to an occurrence of a variable. It passes \emph{through} an
abstraction and chooses left or right at an application \book{187}. Paths contain binders too:
\begin{lprolog}
kind path type.
type bnd          (path -> path) -> path.
type left, right  path -> path.
type path         tm -> path -> o.
path (app M _) (left P) &
path (app _ M) (right P) :- path M P.
path (abs R) (bnd S) :- pi x\ pi p\ path x p => path (R x) (S p).
\end{lprolog}
The \lp{abs} clause introduces \emph{two} eigenvariables at once: a term variable \lp{x} and a path variable
\lp{p}. The assumption \lp{path x p} links them: the path to the occurrence of \lp{x} \emph{is} \lp{p}. The paths of
\lp{(abs x\ app x (abs y\ app y x))} are \book{188}
\begin{lprolog}
bnd u\ left u
bnd u\ right (bnd v\ left v)
bnd u\ right (bnd v\ right u)
\end{lprolog}
The relation also runs backwards: from a set of paths it synthesises the term, using \lp{foreach} from
\cref{ch:hoprog} (\lp{foreach P nil. foreach P (X::L) :- P X, foreach P L.}) \book{188}:
\begin{lprolog}
?- foreach (path N)
     ((bnd u\ left u) ::
      (bnd u\ right (bnd v\ left v))::
      (bnd u\ right (bnd v\ right u))::nil).
N = abs W1\ app W1 (abs W2\ app W2 W1);
no
\end{lprolog}

\begin{tutorial}[Synthesising a term from its paths]
Take the first path. \lp{path N (bnd u\ left u)} can only use the \lp{abs} clause, so \lp{N = abs R} and
\lp{S = u\ left u}. Under eigenvariables \lp{x}, \lp{p} with \lp{path x p} assumed, the goal is
\lp{path (R x) (left p)}. The \lp{left} clause requires \lp{(R x) = app M _}. Unification solves this pattern
problem with \lp{R := x\ app (M' x) (K x)}, two new logic variables applied to \lp{x}. Then \lp{path (M' x) p}
is solved by the assumption \lp{path x p} itself, giving \lp{M' := x\ x}. The second and third paths fill in
the right branch in the same way. Nothing in the program says \enquote{build an abstraction}. The abstraction
in the answer is produced by unification, which abstracts the eigenvariable \lp{x} out of a term that contains
it. This is the operation that Scheme closures could not perform.
\end{tutorial}

Paths can also cross a \emph{marked} $\beta$-redex. A path through \lp{(app (abs R) N)} goes through the body
\lp{R}; when it reaches the bound variable of \lp{R}, it continues into \lp{N}. Marked redexes are written
\lp{(beta N R)} (Figure~7.5, \book{189}):
\begin{lprolog}
type beta     tm -> (tm -> tm) -> tm.
type addbeta  tm -> tm -> o.
type bpath    tm -> path -> o.

bpath (app M _) (left P)  &
bpath (app _ M) (right P) :- bpath M P.
bpath (abs R)   (bnd S)   :- pi x\ pi p\ bpath x p =>
                                         bpath (R x) (S p).
bpath (beta N R) P :-
   pi x\ (pi Q\ bpath x Q :- bpath N Q) => bpath (R x) P.

addbeta (app (abs R) N) (beta M S) :- addbeta (abs R) (abs S),
                                      addbeta N M.
addbeta (app (app M N) P) (app O Q) :- addbeta (app M N) O,
                                       addbeta P Q.
addbeta (abs R) (abs S) :-
  pi x\ (pi M\ pi N\ addbeta (app x M) (app x N) :- addbeta M N) =>
        (addbeta x x) => addbeta (R x) (S x).
\end{lprolog}
Two features are new.
\begin{itemize}[nosep]
  \item The \lp{beta} clause assumes a \emph{rule}, not a fact: \lp{pi Q\ bpath x Q :- bpath N Q}, \enquote{the
        paths of \lp{x} are the paths of \lp{N}}. The assumed clause has a body and its own universally
        quantified variable \lp{Q}, renamed at each use. It also mentions \lp{N}, a variable of the
        enclosing clause, which it shares rather than renames.
  \item \lp{addbeta} marks every redex. It dispatches on the head of an application: an \lp{abs} head is
        marked, an \lp{app} head recurses, and a bound-variable head \lp{x} is handled by the rule assumed
        when \lp{x}'s binder was crossed. A bare variable is copied by the assumed fact \lp{addbeta x x}.
\end{itemize}
Normalisation then has two stages: collect the \lp{bpath}s, then synthesise with \lp{path} \book{189}:
\begin{lprolog}
?- sigma B\ addbeta (app (abs x\x) (abs x\x)) B, bpath B Path.
Path = bnd W1\ W1;
no
?- foreach (P\ path T P) (bnd (W1\ W1) :: nil).
T = abs W1\ W1;
no
\end{lprolog}
For $K\,(S\,K)$ with $K=\lambda x\lambda y\,x$ and $S=\lambda x\lambda y\lambda z\,((x\,z)\,(y\,z))$
\book{189}--\book{190}, the query
\begin{lprolog}
?- sigma K\ sigma S\ sigma B\
     K = (abs x\ abs y\ x),
     S = (abs x\ abs y\ abs z\ app (app x z) (app y z)),
     addbeta (app K (app S K)) B, bpath B Path.
\end{lprolog}
yields four paths, and synthesis gives
\begin{lprolog}
?- foreach (path T)
     ((bnd W1\ bnd W2\ bnd W3\ left (left (bnd W4\ bnd W5\ W4)))::
      (bnd W1\ bnd W2\ bnd W3\ left (right W3))::
      (bnd W1\ bnd W2\ bnd W3\ right (left W2))::
      (bnd W1\ bnd W2\ bnd W3\ right (right W3))::nil).
T = abs W1\ abs W2\ abs W3\
      app (app (abs W4\ abs W5\ W4) W3) (app W2 W3);
no
\end{lprolog}
The result $\lambda w_1\lambda w_2\lambda w_3\,(((\lambda w_4\lambda w_5\,w_4)\,w_3)\,(w_2\,w_3))$ is \emph{not}
$\beta$-normal. The redex $(K\,w_3)$ was \emph{created} by the reduction, so it was not marked. Iterating
\enquote{mark, collect, synthesise} reaches the normal form if it terminates, and in the untyped calculus it
need not terminate \book{190}.

\begin{pitfall}[Collecting all answers is outside the logic]
The two-stage method collects all solutions of \lp{bpath B Path} into a list. The book stresses that this
\enquote{is not formally possible in the simple proof-theoretic framework that we have described here}
\book{190}. Some Prolog systems provide \lp{bagof}, and persistent storage can simulate it. A miniKanren has
the same situation. A goal that consumes a whole sub-search, such as the aggregation goals of \micro{}
(\cref{ch:micro}), is available, but it is extra-logical. It also terminates only when the sub-search is
finite.
\end{pitfall}

\subsection{Type inference}\label{sec:lambdatree-typeof}

The book's type inference for the untyped calculus (\book{190}--\book{191}) is the best-known
instance of the idiom:
\begin{lprolog}
kind ty type.
type arr     ty -> ty -> ty.
type typeof  tm -> ty -> o.
typeof (app M N) A :- typeof M (arr B A), typeof N B.
typeof (abs M) (arr A B) :- pi x\ typeof x A => typeof (M x) B.
\end{lprolog}
Read as inference rules, the typing context $\Gamma$ is the set of assumed atoms \lp{typeof c A}:
\begin{mathpar}
\inferrule{\Gamma\vdash M:B\to A \\ \Gamma\vdash N:B}{\Gamma\vdash \cst{app}\,M\,N : A}
\and
\inferrule{\Gamma,\ c:A\vdash M\,c:B \\ c \text{ fresh}}{\Gamma\vdash\cst{abs}\,M : A\to B}
\end{mathpar}
The context is not a data structure of the program. It is the set of clauses that the \lp{=>} goals have added,
and the variable rule of the type system is \rulename{backchain} on one of them. The interaction is
\book{191}:
\begin{lprolog}
?- typeof (abs x\ abs y\ abs z\ app (app x z) (app y z)) Ty.
Ty = arr (arr T1 (arr T2 T3)) (arr (arr T1 T2) (arr T1 T3))
?- typeof (abs x\x) Ty.
Ty = arr T1 T1
?- typeof (abs x\ app x x) Ty.
no
\end{lprolog}

\begin{tutorial}[Why \texttt{typeof (abs x\textbackslash x) Ty} answers \texttt{arr T1 T1}]
\begin{enumerate}[nosep]
  \item Backchaining on the \lp{abs} clause binds \lp{Ty = arr A B} for fresh logic variables \lp{A}, \lp{B}, and
        \lp{M = x\x}.
  \item \rulename{generic} creates an eigenvariable \lp{c}. \rulename{augment} adds the clause \lp{typeof c A}.
        This assumed clause contains the \emph{logic variable} \lp{A}, which is shared with the rest of the
        derivation. Assumptions are not closed facts.
  \item The goal is now \lp{typeof ((x\x) c) B}, that is, \lp{typeof c B}. Neither static clause applies, since
        \lp{c} is neither an \lp{app} nor an \lp{abs}. The assumed clause applies and unifies \lp{B} with \lp{A}.
  \item So \lp{Ty = arr A A}, printed \lp{arr T1 T1}.
\end{enumerate}
For \lp{abs x\ app x x} the same steps give the assumption \lp{typeof c A}, and \lp{typeof (app c c) B} needs
\lp{typeof c (arr B0 B)} and \lp{typeof c B0}. With the assumption these are \lp{A = arr B0 B} and \lp{A = B0},
hence \lp{B0 = arr B0 B}. The occurs check rejects this, and the answer is \lp{no}. Without an occurs check the
answer would be a cyclic type.
\end{tutorial}

Two more details matter. The type \lp{ty} has only the constructor \lp{arr}, so it has \emph{no closed
inhabitants}. Inferred types are therefore necessarily open, which links back to the distinction between
unifiers and solutions of \cref{ch:lambda}. With an added constant \lp{type i ty.} one can ask closed
questions \book{191}:
\begin{lprolog}
?- typeof (abs x\x) (arr i i).
yes
?- typeof (abs x\x) (arr i Ty).
Ty = i
\end{lprolog}

\begin{theorem}[Subject reduction for call-by-name, \book{191}--\book{192}]\label{thm:lambdatree-subject}
If \lp{(cbn P V)} and \lp{(typeof P A)} are derivable, then \lp{(typeof V A)} is derivable.
\end{theorem}
\begin{proof}[Proof (book's argument)]
By induction on the derivation of \lp{cbn P V}. If it uses the first \lp{cbn} clause, $P=V$. Otherwise
\lp{P = (app M N)}, and there are shorter derivations of \lp{cbn M (abs R)} and \lp{cbn (R N) V}. From
\lp{typeof (app M N) A} there is a \lp{B} with \lp{typeof M (arr B A)} and \lp{typeof N B}. By the induction
hypothesis, \lp{typeof (abs R) (arr B A)}. Hence \lp{pi x\ typeof x B => typeof (R x) A} is derivable.
Instantiating \lp{x} with \lp{N} and using \lp{typeof N B} gives \lp{typeof (R N) A}. The induction hypothesis on
\lp{cbn (R N) V} gives \lp{typeof V A}.
\end{proof}
The book adds that the same holds for \lp{cbv} (left to the reader), and comments: \enquote{The hard technical
part in such a proof is establishing a substitution lemma \ldots{} Here, that substitution lemma is obtained
essentially for free by referring to well-known properties of intuitionistic logic} \book{192}.

\begin{kanren}[The usual miniKanren rendering: an explicit context]
The standard miniKanren way to write a type inferencer avoids both \lp{pi} and \lp{=>}. It uses de Bruijn
indices, so that bound names do not exist, and it passes the context as an explicit list. This runs on today's
\micro{} (indices are Peano numerals and start from \scm{z} for the innermost binder; the book's de Bruijn
indices in \cref{sec:lambdatree-debruijn} start from 1):
\begin{scheme}
(define-relation (nth° Γ k a)
  (fresh° (b Γ* k*)
    (=° Γ `(,b . ,Γ*))
    (cond° ((=° k 'z) (=° a b))
           ((=° k `(s ,k*)) (nth° Γ* k* a)))))

(define-relation (typeof° Γ m a)
  (cond°
    ((fresh° (k) (=° m `(var ,k)) (nth° Γ k a)))
    ((fresh° (m1 n b) (=° m `(app ,m1 ,n)) (typeof° Γ m1 `(arr ,b ,a)) (typeof° Γ n b)))
    ((fresh° (body a1 b) (=° m `(lam ,body)) (=° a `(arr ,a1 ,b)) (typeof° `(,a1 . ,Γ) body b)))))

(define S '(lam (lam (lam (app (app (var (s (s z))) (var z)) (app (var (s z)) (var z)))))))
(define I '(lam (var z)))
(define ω '(lam (app (var z) (var z))))

(print (μkanren-run (q -1 #t) (typeof° '() S q)))
(print (μkanren-run (q -1 #t) (typeof° '() I q)))
(print (μkanren-run (q -1 #t) (typeof° '() ω q)))
\end{scheme}
\begin{answers}
((arr (arr α (arr β γ)) (arr (arr α β) (arr α γ))))
((arr α α))
()
\end{answers}
These are the book's three answers. The occurs check of \micro{} rejects $\lambda x\,(x\,x)$ exactly as in the
tutorial above. The difference between the two programs is where the context lives. In \scm{typeof°} it is
data, threaded through every call and looked up with a hand-written \scm{nth°}. In \lp{typeof} it is part of
the program, and the logic guarantees the substitution lemma behind \cref{thm:lambdatree-subject}. With an
explicit context, every relation that may meet a bound variable must accept and pass on the context. With
\lp{=>}, a predicate written without any knowledge of binders, such as \lp{atom} in
\cref{sec:lambdatree-fol-comp}, keeps working under them. Supporting the second style is the point of the
proposal (\cref{ch:design}).
\end{kanren}

\subsection{Translating to and from de Bruijn syntax}\label{sec:lambdatree-debruijn}

In de Bruijn's notation \cite{deBruijn1972} a variable is a positive integer $n$ that refers to the $n$-th
enclosing abstraction, counting from 1 \book{192}:
\begin{center}
\begin{tabular}{@{}ll@{\qquad}ll@{}}
\toprule
term & de Bruijn & term & de Bruijn\\
\midrule
$\lambda x\,x$ & $\lambda 1$ & $\lambda x\lambda y\,(y\,x)$ & $\lambda\lambda(1\,2)$\\
$\lambda x\,(x\,x)$ & $\lambda(1\,1)$ & $\lambda x\lambda y\lambda z\,((x\,z)(y\,z))$ & $\lambda\lambda\lambda((3\,1)(2\,1))$\\
$\lambda x\lambda y\,x$ & $\lambda\lambda2$ & $\lambda x\,(x\,(\lambda y\,x))$ & $\lambda(1\,\lambda2)$\\
$\lambda x\lambda y\,y$ & $\lambda\lambda1$ & $\lambda x\,(x\,(\lambda y\,(x\,(\lambda w\,(w\,x)))))$ & $\lambda(1\,(\lambda(2\,(\lambda(1\,3)))))$\\
\bottomrule
\end{tabular}
\end{center}
The same variable receives different numbers at different depths (the $x$ of the last row is 1, 2 and 3).
The encoding \book{193}, in which \lp{deb} is both a kind and a constructor, and the translation are:
\begin{lprolog}
kind deb  type.
type ab   deb -> deb.
type ap   deb -> deb -> deb.
type deb  int -> deb.
\end{lprolog}
\begin{lprolog}
type trans  int -> tm -> deb -> o.
type depth  int -> tm -> o.
trans D (abs M) (ab P) :- pi c\ depth D c =>
                            (E is D + 1, trans E (M c) P).
trans D (app M N) (ap P Q) :- trans D M P, trans D N Q.
trans D X (deb E) :- depth N X, E is (D - N).
\end{lprolog}
In \lp{trans D t d}, \lp{D} is one more than the number of enclosing binders. Each eigenvariable is tagged, by
the assumption \lp{depth D c}, with the level at which it was bound. At an occurrence, the index is the current
level minus the binding level. The relation works in both directions \book{193}:
\begin{lprolog}
?- trans 1 (abs x\ app x (abs y\ app x (abs w\ app w x))) D.
D = ab (ap (deb 1) (ab (ap (deb 2) (ab (ap (deb 1) (deb 3))))))
?- trans 1 P
     (ab (ap (deb 1) (ab (ap (deb 2) (ab (ap (deb 1) (deb 3))))))).
P = abs W1\ app W1 (abs W2\ app W1 (abs W3\ app W3 W1))
\end{lprolog}

\begin{tutorial}[Tracing \texttt{?- trans 1 (abs x\textbackslash\ abs y\textbackslash\ abs z\textbackslash\ y) P.}]
The book traces this query \book{193}--\book{194}.
\begin{enumerate}[nosep]
  \item \lp{P := (ab P1)}. A new eigenvariable \lp{c} is introduced and \lp{depth 1 c} is assumed. Goal:
        \lp{trans 2 (abs y\ abs z\ y) P1}. (The meta-level $\beta$ has replaced the bound \lp{x} by \lp{c}; since
        \lp{x} does not occur in the body, nothing visible changes.)
  \item \lp{P1 := (ab P2)}, eigenvariable \lp{d}, assume \lp{depth 2 d}. Goal: \lp{trans 3 (abs z\ d) P2}.
  \item \lp{P2 := (ab P3)}, eigenvariable \lp{e}, assume \lp{depth 3 e}. Goal: \lp{trans 4 d P3}.
  \item Only the third clause applies: \lp{P3 := (deb E)} and \lp{depth N d, E is (4 - N)}. The assumption gives
        \lp{N = 2}, so \lp{E = 2}.
  \item Answer: \lp{P = (ab (ab (ab (deb 2))))}.
\end{enumerate}
The book's text writes the assumptions as \enquote{depth c 1}, \enquote{depth d 2}, \enquote{depth e 3}, with the
arguments swapped with respect to \lp{type depth int -> tm -> o}. The program assumes \lp{depth 1 c}, and so on.
In the reverse direction the same clauses run with \lp{P} unknown. The third clause then finds the eigenvariable
\lp{X} whose \lp{depth} gives the right index, and the \lp{(M c)} of the first clause is solved as a pattern
problem, as in the path synthesis of \cref{sec:lambdatree-paths}.
\end{tutorial}

A variant (\book{194}) with no \lp{depth} predicate assumes new \lp{trans} clauses for each
eigenvariable. It has \emph{no} static clause for variables:
\begin{lprolog}
trans D (app M N) (ap P Q) :- trans D M P, trans D N Q.
trans D (abs M) (ab P) :-
   pi u\
     (pi N\ pi H\ trans N u (deb H) :- H is (N - D))
       => (D' is D + 1, trans D' (M u) P).
\end{lprolog}
For the query traced above, the assumed clauses are
\begin{lprolog}
pi N\ pi H\ trans N c (deb H) :- H is (N - 1).
pi N\ pi H\ trans N d (deb H) :- H is (N - 2).
pi N\ pi H\ trans N e (deb H) :- H is (N - 3).
\end{lprolog}
and \lp{?- trans 4 d P3.} is solved by the second. Each assumed clause has captured the value of the outer
\lp{D} at the time it was assumed. It behaves as a closure over the context.

\begin{kanren}[The same translation with an explicit environment]
In \micro{} the translation is written with an explicit environment of names. It is the miniKanren analogue of
the \lp{depth} assumptions:
\begin{scheme}
; named terms: a symbol, (app m n), (λ x b); de Bruijn terms as in typeof.scm (z = innermost binder)
(define-relation (index° env x k)
  (fresh° (y env*)
    (=° env `(,y . ,env*))
    (cond° ((=° x y) (=° k 'z))
           ((fresh° (k*) (≠° x y) (=° k `(s ,k*)) (index° env* x k*))))))

(define-relation (trans° env t d)
  (cond°
    ((symbol° t) (fresh° (k) (=° d `(var ,k)) (index° env t k)))
    ((fresh° (m n p r) (=° t `(app ,m ,n)) (=° d `(app ,p ,r)) (trans° env m p) (trans° env n r)))
    ((fresh° (x b p) (=° t `(λ ,x ,b)) (symbol° x) (=° d `(lam ,p)) (trans° `(,x . ,env) b p)))))

; λx (x (λy (x (λw (w x))))), the p.193 example, forwards
(print (μkanren-run (q -1 #t) (trans° '() '(λ x (app x (λ y (app x (λ w (app w x)))))) q)))
; λx λy x backwards: the binder names come back as constrained logic variables
(print (μkanren-run (q -1 #f) (trans° '() q '(lam (lam (var (s z)))))))
\end{scheme}
\begin{answers}
((lam (app (var z) (lam (app (var (s z)) (lam (app (var z) (var (s (s z))))))))))
((λ (α β) (begin (deny (equal? β α))) (assert (every (μ v (symbol? v)) (list α β))) (cons (quote λ) (cons α (cons (cons (quote λ) (cons β (cons α (quote ())))) (quote ()))))))
\end{answers}
The forward answer is the book's \lp{D}, shifted to 0-based indices. The backward answer shows what names cost.
λProlog answers \lp{abs W1\ abs W2\ W1}: there are no names, and the printer invents \lp{W1}, \lp{W2}. The named
miniKanren answer is $(\lambda\,\alpha\,(\lambda\,\beta\,\alpha))$ for \emph{any} two symbols $\alpha\ne\beta$.
The disequality is needed to rule out the shadowed reading $\lambda x\lambda x\,x$, which would be
$\lambda\lambda1$. The names are logic variables with side conditions, and the programmer had to write the
\scm{≠°} that excludes capture. In λ-tree syntax there is no such condition to write.
\end{kanren}

\section{Computations over first-order formulas}\label{sec:lambdatree-fol-comp}

\subsection{Vacuous quantifiers and quantifier-free formulas}

Figure~7.6 (\book{195}):
\begin{lprolog}
type vacuous    form -> form -> o.
type quantfree  form -> o.
vacuous (all  x\ P) P.
vacuous (some x\ P) P.
quantfree tt & quantfree ff.
quantfree F :- atom F.
quantfree (neg F) :- quantfree F.
quantfree (F && G) & quantfree (F !! G) & quantfree (F ==> G) :-
   quantfree F, quantfree G.
\end{lprolog}
\lp{vacuous} \enquote{succeeds exactly when the top-level abstracted variable does not appear in the body of the
formula and, in this case, binds the second argument to the body} \book{195}. There is no occurrence test in
the code. Matching \lp{(all x\ P)} against \lp{(all x\ p a)} succeeds with \lp{P = p a}. Matching it against
\lp{(all x\ p x)} would need \lp{P = p x}, but \lp{P} is quantified outside \lp{x\}, so \lp{x} cannot appear in its
value and unification fails. The scoping constraint performs the occurrence test. \lp{quantfree} relies on
\lp{atom} (as in Figure~7.2), which is defined for each object signature.

\subsection{Negation normal form}\label{sec:lambdatree-nnf}

A formula is in \emph{negation normal form} (NNF) if it contains no implication and negation is applied only to
atoms. It is in \emph{prenex normal form} if no quantifier is in the scope of a connective. In classical
first-order logic every formula is equivalent to one in NNF, to one in prenex form, and to one in both
\book{195}. The NNF transformation orients the equivalences
\[
\begin{aligned}
\neg\neg B&\equiv B, & B_1\imp B_2&\equiv\neg B_1\vee B_2, & \neg\forall x\,B\,x&\equiv\exists x\,\neg B\,x,\\
\neg(B_1\wedge B_2)&\equiv\neg B_1\vee\neg B_2, & \neg(B_1\vee B_2)&\equiv\neg B_1\wedge\neg B_2, &
\neg\exists x\,B\,x&\equiv\forall x\,\neg B\,x
\end{aligned}
\]
from left to right (Figure~7.7, \book{196}):
\begin{lprolog}
type nnf  form -> form -> o.
nnf A A & nnf (neg A) (neg A) :- atom A.
nnf (neg (neg B)) D :- nnf B D.
nnf (neg (B && C)) (D !! E) &
nnf (neg (B !! C)) (D && E) :- nnf (neg B) D, nnf (neg C) E.
nnf (B ==> C) (D !! E) :- nnf (neg B) D, nnf C E.
nnf (B !! C) (D !! E) &
nnf (B && C) (D && E) :- nnf B D, nnf C E.
nnf (neg (all B)) (some D) &
nnf (neg (some B)) (all D) :- pi x\ nnf (neg (B x)) (D x).
nnf (all B) (all D) &
nnf (some B) (some D) :- pi x\ nnf (B x) (D x).
\end{lprolog}
The quantifier clauses are the idiom without \lp{=>}. Under \lp{pi x\} the eigenvariable \lp{x} is an object
\emph{term}, and the only predicate that meets the atoms containing it is \lp{atom}, which does not look inside
terms. The output \lp{(D x)} is again a pattern problem. Its solution abstracts \lp{x} out of the transformed
body, rebuilding the quantifier around it. For example (our trace, assuming \lp{atom} holds of \lp{p} atoms),
\lp{?- nnf (neg (all x\ p x)) D.} skips the two \lp{atom}-guarded clauses, since neither the whole formula nor
\lp{all x\ p x} is atomic. It reaches the \lp{neg (all B)} clause with \lp{B = x\ p x} and \lp{D = some D1}. Under
an eigenvariable \lp{c} the goal is \lp{nnf (neg (p c)) (D1 c)}. The clause \lp{nnf (neg A) (neg A) :- atom A}
gives \lp{(D1 c) = neg (p c)}, so \lp{D1 := x\ neg (p x)} and \lp{D = some x\ neg (p x)}, the formula
$\exists x\,\neg p(x)$.

The book claims that no two heads overlap, so \lp{nnf} is a partial function, and that \enquote{a simple argument
by induction shows also that the function represented is total} \book{196}. Both claims need qualification.

\begin{pitfall}[Non-overlap modulo a guard, and a missing case]
Syntactically, the head \lp{nnf A A} overlaps every other head, since \lp{A} is a variable. The clauses are
disjoint only \emph{modulo the guard} \lp{atom A}, which fails on compound formulas. A miniKanren port must keep
that guard decisive. The \lp{atom} relation must \emph{fail}, not enumerate or suspend, on compound or unbound
input, or \lp{nnf} produces spurious answers. Moreover, \lp{tt}, \lp{ff} and \lp{neg (B ==> C)} have no clause.
The last matches none of \lp{neg (neg _)}, \lp{neg (_ && _)}, \lp{neg (_ !! _)}, \lp{neg (all _)},
\lp{neg (some _)}. Totality holds only for the formulas without these cases. The fix when porting is one more
clause, \lp{nnf (neg (B ==> C)) (D && E) :- nnf B D, nnf (neg C) E.}
\end{pitfall}

\subsection{Prenex normal form}\label{sec:lambdatree-prenex}

For formulas in NNF, the prenex transformation uses the equivalences \book{196}
\[
(\forall x\,B_1x)\wedge(\forall x\,B_2x)\equiv\forall x\,(B_1x\wedge B_2x),\qquad
(\exists x\,B_1x)\vee(\exists x\,B_2x)\equiv\exists x\,(B_1x\vee B_2x),
\]
and, for $Q\in\{\forall,\exists\}$ and $\circ\in\{\wedge,\vee\}$,
$B_1\circ(Q x\,B_2x)\equiv Qx\,(B_1\circ B_2x)$ and $(Qx\,B_2x)\circ B_1\equiv Qx\,(B_2x\circ B_1)$. The usual
side condition, $x$ not free in $B_1$, is not written anywhere in the program, because it is automatic: in the
clauses below \lp{B} and \lp{C} are clause variables quantified outside the \lp{pi x\}. (The restriction to NNF
input \enquote{is not necessary and is used only to shorten the specification}, \book{196}.) Figure~7.8
(\book{197}):
\begin{lprolog}
type prenex, merge  form -> form -> o.
prenex A A & prenex (neg A) (neg A) :- atom A.
prenex (B && C) D :- prenex B U, prenex C V,
                     merge (U && V) D.
prenex (B !! C) D :- prenex B U, prenex C V,
                     merge (U !! V) D.
prenex (all B) (all D) &
prenex (some B) (some D) :- pi x\ prenex (B x) (D x).
merge ((all B) && (all C)) (all D) :-
   pi x\ merge ((B x) && (C x)) (D x).
merge ((all B) && C) (all D) &
merge ((some B) && C) (some D) :-
   pi x\ merge ((B x) && C) (D x).
merge (B && (all C)) (all D) &
merge (B && (some C)) (some D) :-
   pi x\ merge (B && (C x)) (D x).
merge ((some B) !! (some C)) (some D) :-
   pi x\ merge ((B x) !! (C x)) (D x).
merge ((some B) !! C) (some D) &
merge ((all B) !! C) (all D) :-
   pi x\ merge ((B x) !! C) (D x).
merge (B !! (some C)) (some D) &
merge (B !! (all C)) (all D) :-
   pi x\ merge (B !! (C x)) (D x).
merge B B :- quantfree B.
\end{lprolog}
Look at the clause \lp{merge ((all B) && C) (all D) :- pi x\ merge ((B x) && C) (D x).} The formula \lp{C}
moves under the new quantifier without any renaming, because \lp{C} is a closed value with respect to \lp{x}.
Capture is impossible by construction.

Some formulas have a unique prenex form \book{196}. For example
\begin{lprolog}
(all x\ (p x) && (all y\ q x y) && (p (f x))) !! (p a).
\end{lprolog}
has only
\begin{lprolog}
all x\ all y\ ((p x) && (q x y) && (p (f x))) !! (p a).
\end{lprolog}
In general, though, \lp{prenex} is a relation, not a function \book{196}--\book{197}:
\begin{lprolog}
?- prenex ((all x\ q x x) && (all z\ all y\ q z y)) P.
\end{lprolog}
has five answers \book{197}:
\begin{lprolog}
all z\ all y\ (q z z) && (q z y)
all x\ all z\ all y\ (q x x) && (q z y)
all z\ all x\ (q x x) && (q z x)
all z\ all x\ all y\ (q x x) && (q z y)
all z\ all y\ all x\ (q x x) && (q z y)
\end{lprolog}

\begin{tutorial}[Where the five prenex forms come from]
Both conjuncts are already prenex, so the goal reduces to
\lp{merge ((all x\ q x x) && (all z\ all y\ q z y)) P}. The following explanation is our own hand trace
against the clause order of Figure~7.8. It reproduces the book's order under depth-first search.
\begin{enumerate}[nosep]
  \item The $\forall\wedge\forall$ clause shares one eigenvariable between \lp{x} and \lp{z}, leaving
        \lp{merge ((q c c) && (all y\ q c y))}, which pulls \lp{y} out: answer 1.
  \item The \lp{((all B) && C)} clause pulls \lp{x} alone, then \lp{z} and \lp{y} come out of the right side one at
        a time: answer 2.
  \item The \lp{(B && (all C))} clause pulls \lp{z} first, leaving
        \lp{merge ((all x\ q x x) && (all y\ q c y))}. That subproblem again has three ways out: the
        $\forall\wedge\forall$ clause sharing \lp{x} with \lp{y} (answer 3), pulling \lp{x} then \lp{y} (answer
        4), or pulling \lp{y} then \lp{x} (answer 5).
\end{enumerate}
Every answer is a correct prenex form. No other clause applies at any step: the $\vee$ clauses need \lp{!!},
and \lp{merge B B :- quantfree B} fails while a quantifier remains. A miniKanren port with the same clauses
would have to produce exactly these five answers, but its interleaving search may list them in a different
order.
\end{tutorial}

\subsection{Syntactic classes and a richer interpreter}

Figure~7.9 (\book{198}) recognises the goal and clause formulas of \fohc{} and \fohh{} (\cref{ch:horn,ch:hh}):
\begin{lprolog}
type fohcG, fohcD, fohhG, fohhD  form -> o.
fohcG tt.
fohcG A :- atom A.
fohcG (some B) :- pi x\ fohcG (B x).
fohcG (B && C) & fohcG (B !! C) :- fohcG B, fohcG C.
fohcD A :- atom A.
fohcD (G ==> D) :- fohcG G, fohcD D.
fohcD (D1 && D2) :- fohcD D1, fohcD D2.
fohcD (all D) :- pi x\ fohcD (D x).
fohhG tt.
fohhG A :- atom A.
fohhG (D ==> G) :- fohhD D, fohhG G.
fohhG (B && C) & fohhG (B !! C) :- fohhG B, fohhG C.
fohhG (some B) & fohhG (all B) :- pi x\ fohhG (B x).
fohhD A :- atom A.
fohhD (D1 && D2) :- fohhD D1, fohhD D2.
fohhD (G ==> D) :- fohhG G, fohhD D.
fohhD (all D) :- pi x\ fohhD (D x).
\end{lprolog}
Here the book remarks on the \lp{pi x\} without \lp{=>}, a deviation from the idiom \book{198}. The
eigenvariables are object terms and \lp{atom} does not inspect terms, so no assumptions are needed. (The same
silent deviation occurs in \lp{red1}, \lp{nnf} and \lp{prenex}.) The book also notes a free substitution theorem:
\enquote{If $\forall x\,D$ is a first-order Horn clause and t is a term, then $D[t/x]$ is a first-order Horn
clause} \book{198}. It is the analogue of \cref{thm:lambdatree-bnorm-subst} and is proved the same way.

Two clauses extend the interpreter of Figure~7.1 from \fohc{} to \fohh{} \book{198}:
\begin{lprolog}
interp D (D1 ==> G) :- interp (D1 && D) G.
interp D (all G) :- pi x\ interp D (G x).
\end{lprolog}
Object $\wedge$, $\vee$, $\exists$ in goals are handled by meta $\wedge$, $\vee$, $\exists$, and the object
$\forall$ in goals by the meta \lp{pi}. Object implication in goals extends the object program. Since the object
program is an explicit argument, the extension is just \lp{(D1 && D)}, and no meta-level \lp{=>} is used.

\section{Specifying substitution without \texorpdfstring{$\beta$}{β}}\label{sec:lambdatree-copy}

Substitution into a term with an abstraction, by $\beta$, is one clause \book{199}:
\begin{lprolog}
type subst (tm -> tm) -> tm -> tm -> o.
subst R N (R N).
\end{lprolog}
As \cref{ch:lambda} showed, however, \enquote{full $\beta$-conversion is a complex operation on terms} \book{199}. The book therefore shows how to
specify substitution with only the weak $\beta_0$ of \cref{sec:lambdatree-llambda}. The tools are a
\emph{copy} predicate and a hypothetical assumption.

\subsection{Copy clauses}

Figure~7.10 (\book{200}):
\begin{lprolog}
type copy   tm -> tm -> o.
type subst  (tm -> tm) -> tm -> tm -> o.
copy (app M N) (app P Q) :- copy M P, copy N Q.
copy (abs M) (abs N) :- pi x\ copy x x => copy (M x) (N x).
subst M T S :- pi x\ copy x T => copy (M x) S.
\end{lprolog}
\lp{copy} specifies equality: for closed terms, \lp{copy t s} is derivable if and only if $t=s$. A trace
\book{199}:
\begin{lprolog}
?- copy (abs x\ abs y\ app y x) M.
   % M := (abs M1); eigenvariable c; assume (copy c c)
?- copy (abs y\ app y c) (M1 c).
   % (M1 c) := (abs M2); eigenvariable d; assume (copy d d)
?- copy (app d c) (M2 d).
   % three backchainings: app clause, then (copy d d), then (copy c c)
   % finally (M2 d) unified with (app d c)
M2 = u\app u c
M1 = v\abs u\app u v
M  = abs v\abs u\app u v.
\end{lprolog}
The equations \lp{(M1 c) = (abs M2)} and \lp{(M2 d) = (app d c)} are pattern problems, solved by abstracting
the eigenvariable out. So \lp{copy} \emph{rebuilds} its input binder by binder, and the abstractions of the output
are produced by unification.

\begin{tutorial}[Reading \texttt{subst} as a hypothetical copy]
\lp{subst M T S} copies the body of \lp{M} except that the outermost bound variable copies to \lp{T} \book{200}.
The clause does this in three steps:
\begin{enumerate}[nosep]
  \item the binder of \lp{M} becomes a formula-level \lp{pi x\};
  \item the assumption \lp{copy x T} says \enquote{wherever you meet \lp{x}, produce \lp{T}};
  \item \lp{copy (M x) S} copies the body.
\end{enumerate}
Inside the body, the assumptions \lp{copy c c} for inner binders preserve inner variables. The assumption
\lp{copy x T} replaces the substituted one. Since \lp{T} is closed with respect to every eigenvariable created
inside, no capture can happen.
\end{tutorial}

\begin{proposition}[Correctness of \texttt{subst}, \book{200}]\label{prop:lambdatree-subst}
If \lp{(subst M T S)} is derivable, then \lp{S} is equal to \lp{(M T)}.
\end{proposition}
\begin{proof}
\lp{pi x\ copy x T => copy (M x) S} is derivable. Since \lp{copy} is equality in the program without
assumptions, \lp{copy T T} is derivable. Instantiate \lp{x} with \lp{T}: the assumption becomes \lp{copy T T} and can
be discharged by modus ponens. So \lp{copy (M T) S} is derivable from the program without assumptions, hence
\lp{(M T)} $=$ \lp{S}.
\end{proof}

\subsection{Copy clauses for first-order formulas, and where they come from}

For the first-order formulas over the signature
\begin{lprolog}
type a  term.
type f  term -> term.
type g  term -> term -> term.
type p  term -> form.
type q  term -> term -> form.
\end{lprolog}
the copy clauses and substitution are Figure~7.11 (\book{201}):
\begin{lprolog}
kind term, form  type.
type copyterm  term -> term -> o.
type copyform  form -> form -> o.
type subst     (term -> form) -> term -> form -> o.

copyterm a a.
copyterm (f X) (f U) :- copyterm X U.
copyterm (g X Y) (g U V) :- copyterm X U, copyterm Y V.

copyform tt tt.
copyform ff ff.
copyform (neg B) (neg D) :- copyform B D.
copyform (B && C) (D && E) &
copyform (B !! C) (D !! E) &
copyform (B ==> C) (D ==> E) :- copyform B D, copyform C E.
copyform (all B) (all D) &
copyform (some B) (some D) :- pi y\ copyterm y y =>
                                     copyform (B y) (D y).
copyform (p X) (p U) :- copyterm X U.
copyform (q X Y) (q U V) :- copyterm X U, copyterm Y V.

subst M T N :- pi x\ copyterm x T => copyform (M x) N.
\end{lprolog}
These clauses are not ad hoc. They are generated from the signature by a translation that \enquote{lowers} the
equality at a higher type to relations at the base types \book{201}--\book{202}.

\begin{definition}[Type lowering, \book{201}--\book{202}]\label{def:lambdatree-lowering}
For terms $t,s$ of type $\tau$ built from the base types \lp{term} and \lp{form}, define the formulas
$[\![t,s:\tau]\!]^{+}$ (for goals) and $[\![t,s:\tau]\!]^{-}$ (for clauses) by
\[
\begin{aligned}
{}[\![t,s:\cst{term}]\!]^{\pm} &= \cst{copyterm}\ t\ s, \qquad
[\![t,s:\cst{form}]\!]^{\pm} = \cst{copyform}\ t\ s,\\
[\![t,s:\tau\to\sigma]\!]^{+} &= \forall x\,\big([\![x,x:\tau]\!]^{-}\imp[\![t\,x,\ s\,x:\sigma]\!]^{+}\big),\\
[\![t,s:\tau\to\sigma]\!]^{-} &= \forall x\,\forall y\,\big([\![x,y:\tau]\!]^{+}\imp[\![t\,x,\ s\,y:\sigma]\!]^{-}\big).
\end{aligned}
\]
\end{definition}
The positive case is the idiom itself: to compare two functions, apply both to a fresh $x$ and assume that $x$
is related to itself. The negative case, used for clauses, compares arguments pairwise. Polarity flips at each
arrow, just as an arrow-typed argument of a constructor turns a clause into a hypothetical goal.

\begin{example}[Lowering the quantifier \lp{all}, \book{201}--\book{202}]\label{ex:lambdatree-lower-all}
For $\cst{all}:(\cst{term}\to\cst{form})\to\cst{form}$:
\[
\begin{aligned}
&[\![\cst{all},\cst{all}:(\cst{term}\to\cst{form})\to\cst{form}]\!]^{-}\\
&\quad= \forall x\forall y\,\big([\![x,y:\cst{term}\to\cst{form}]\!]^{+}\imp[\![\cst{all}\,x,\cst{all}\,y:\cst{form}]\!]^{-}\big)\\
&\quad= \forall x\forall y\,\big((\forall z\,[\![z,z:\cst{term}]\!]^{-}\imp[\![x\,z,\ y\,z:\cst{form}]\!]^{+})\imp[\![\cst{all}\,x,\cst{all}\,y:\cst{form}]\!]^{-}\big)\\
&\quad= \forall x\forall y\,\big(\forall z\,(\cst{copyterm}\ z\ z\imp\cst{copyform}\ (x\,z)\ (y\,z))\\
&\qquad\qquad\imp\cst{copyform}\ (\cst{all}\,x)\ (\cst{all}\,y)\big).
\end{aligned}
\]
This is the \lp{all} clause of Figure~7.11. \enquote{In fact, the clauses in that figure are exactly the formulas
$[\![c, c : \tau]\!]^-$ for each of the constants c : $\tau$ of the object logic} \book{202}.
\end{example}

\begin{proposition}[\book{202}]
\lp{(copyterm T S)} is derivable if and only if \lp{T} and \lp{S} are equal terms, and \lp{(copyform B C)} is
derivable if and only if \lp{B} and \lp{C} are equal formulas. Hence, if \lp{(subst M T S)} is provable, then
\lp{S} is equal to \lp{(M T)}.
\end{proposition}

Two remarks from \book{202}. First, these programs are \emph{signature-dependent}: the object signature is
reflected into clauses, whereas the one-line $\beta$ version is not. Second, lowering is general. Each primitive
type maps to its own predicate, the relation need not be binary, and the \lp{term} predicate of
\cref{sec:lambdatree-term} is the arity-one lowering for the signature \lp{abs}, \lp{app}.

Substitution can also be defined by direct recursion on the abstraction, with no copy predicate (Figure~7.12,
\book{203}):
\begin{lprolog}
type subst      (term -> form) -> term -> form -> o.
type substterm  (term -> term) -> term -> term -> o.
subst (x\ tt) T tt.
subst (x\ ff) T ff.
subst (x\ neg (B x)) T (neg D) :- subst B T D.
subst (x\ (B x) && (C x))  T (D && E)  &
subst (x\ (B x) !! (C x))  T (D !! E)  &
subst (x\ (B x) ==> (C x)) T (D ==> E) :- subst B T D,
                                         subst C T E.
subst (x\ all  (B x)) T (all  D) &
subst (x\ some (B x)) T (some D) :-
   pi y\ substterm (x\y) T y => subst (x\ B x y) T (D y).
subst (x\ p (X x)) T (p U) :- substterm X T U.
subst (x\ q (X x) (Y x)) T (q U V) :- substterm X T U,
                                      substterm Y T V.
substterm (x\ x) T T.
substterm (x\ a) T a.
substterm (x\ f (F x)) T (f S) :- substterm F T S.
substterm (x\ g (F x) (G x)) T (g S R) :-
   substterm F T S,
   substterm G T R.
\end{lprolog}
The book says only that these clauses are signature dependent, and in its notes that they are loosely based
on Andrews' equational presentation of $\lambda$-conversion \book{210}. Three points are worth seeing. The patterns
\lp{(x\ a)} and \lp{(x\y)} are \emph{vacuous} abstractions: they match only bodies that do not depend on \lp{x}.
Under a nested quantifier the body becomes a two-place abstraction \lp{(x\ B x y)}. The assumption
\lp{substterm (x\y) T y} says that the inner bound \lp{y}, seen as a constant function of \lp{x}, substitutes to
itself.

\subsection{Substitution runs backwards}

Every version of \lp{subst} is a relation (\book{202}). Given \lp{T} and the result, it computes
the abstraction:
\begin{lprolog}
?- subst M (f a) ((p (f a)) && (q a (f a))).
M = x\ (p x) && (q a x)
M = x\ (p x) && (q a (f a))
M = x\ (p (f a)) && (q a x)
M = x\ (p (f a)) && (q a (f a))
\end{lprolog}
Each of the two occurrences of \lp{(f a)} may or may not be abstracted, so there are $2\times2$ answers. They
are the four \emph{existential generalisations} of $p(f(a))\wedge q(a,f(a))$ on the term $f(a)$ \book{203}:
$\exists x\,p(x)\wedge q(a,x)$, $\exists x\,p(x)\wedge q(a,f(a))$, $\exists x\,p(f(a))\wedge q(a,x)$,
$\exists x\,p(f(a))\wedge q(a,f(a))$. The \enquote{abstracted} answer comes first. This is consistent with the
convention of \cref{ch:hh} that clauses added by an implication are tried before the static ones, so
\lp{copyterm x (f a)} is tried before \lp{copyterm (f X) (f U)}.

\begin{pitfall}[Copy is equality logically, not operationally]
\enquote{the operational behavior of the copy predicate can be quite different from that of the equality notion
it encodes} \book{207}. \lp{?- M = N.} binds \lp{M} to \lp{N} at once. \lp{?- copy M N.} builds terms. With
\begin{lprolog}
kind i     type.
type a     i.
type f     i -> i.
type copyi i -> i -> o.
copyi a a.
copyi (f X) (f Y) :- copyi X Y.
\end{lprolog}
\lp{?- copyi M N.} enumerates \lp{a}, \lp{(f a)}, \lp{(f (f a))}, \ldots, and \enquote{if the order of the clauses
for copyi were reversed, then this query would cause the depth-first $\lambda$Prolog interpreter to loop without
providing any solution} \book{207}.
\end{pitfall}

\begin{kanren}[Clause order and interleaving]
The same predicate in \micro, in both clause orders:
\begin{scheme}
(define-relation (copyi° m n)
  (cond° ((=° m 'a) (=° n 'a))
         ((fresh° (x y) (=° m `(f ,x)) (=° n `(f ,y)) (copyi° x y)))))
(define-relation (copyi-rev° m n)
  (cond° ((fresh° (x y) (=° m `(f ,x)) (=° n `(f ,y)) (copyi-rev° x y)))
         ((=° m 'a) (=° n 'a))))

(print (μkanren-run (q 4 #t) (fresh° (m n) (=° q `(,m ,n)) (copyi° m n))))
(print (μkanren-run (q 4 #t) (fresh° (m n) (=° q `(,m ,n)) (copyi-rev° m n))))
\end{scheme}
\begin{answers}
((a a) ((f a) (f a)) ((f (f a)) (f (f a))) ((f (f (f a))) (f (f (f a)))))
((a a) ((f a) (f a)) ((f (f a)) (f (f a))) ((f (f (f a))) (f (f (f a)))))
\end{answers}
The reversed order loops in λProlog but produces the same answers here. \scm{define-relation} delays the
recursive call and \scm{cond°} interleaves its branches, so the base case is reached. Interleaving changes the
\emph{order} of answers (the five prenex forms come out in the order 1~3~2~4~5, \cref{sec:examples-prenex};
the four generalisations happen to keep the book's order, \cref{sec:examples-subst}) and turns some depth-first loops into
productive streams. It does not remove genuine divergence. \lp{cbv} on $(\lambda x\lambda w\,w)\,\Omega$ has no
answer under any strategy, and a conjunction such as \enquote{\scm{copyi°} and then fail} still runs forever,
because it has infinitely many candidates and no finite failure.
\end{kanren}

\section{The \texorpdfstring{$\lambda$}{λ}-tree approach to abstract syntax}\label{sec:lambdatree-hoas}

Concrete syntax (strings, whitespace, parentheses, infix operators, comments) is for humans. Parsing produces
\emph{abstract syntax}, a tree \book{203}. For binders there are several choices \book{204}: keep explicit names;
or map object binding to the meta-language's binding, which is called \emph{higher-order abstract syntax}
(HOAS) \cite{PfenningElliott1988}. The term covers \enquote{a range of encoding techniques}, and the book
distinguishes two of them sharply.
\begin{itemize}[nosep]
  \item In functional programming, and in constructive type theories that compute by evaluation, meta-level
        abstraction builds \emph{rich function spaces}. \enquote{Because of the richness of these spaces, it is not
        possible to compare function-valued expressions for equality in any simple way}. If quantified formulas
        are encoded with such functions, \enquote{we cannot also compare or structurally analyze formulas under
        such an encoding} \book{204}.
  \item Most logic programming languages have no meta-level binders in terms at all.
  \item The $\lambda$-terms of hohh use \enquote{a $\lambda$-calculus with a weak notion of functionality that makes
        it possible to analyze the structures of objects under the encoding} \book{204}.
\end{itemize}
The book contrasts \enquote{the two extremes in the treatment of binding---one using functions that are best
thought of extensionally and the other using $\lambda$-expressions that are easily treated intensionally}.

\begin{definition}[$\lambda$-tree syntax, \book{204}]\label{def:lambdatree-lambdatree}
\enquote{We shall use the term $\lambda$-tree syntax to denote the approach to higher-order abstract syntax that uses
typed $\lambda$-terms modulo equality based on $\alpha$-conversion, some sufficiently weak form of
$\beta$-conversion, and possibly $\eta$-conversion.}
\end{definition}
In this chapter $\beta$ is \enquote{weakened} by simple typing, which makes it terminating. The next section
weakens it much further, to $\beta_0$.

The book's bibliographic notes (\book{208}--\book{210}) trace the idea back a long way. Encoding a binder as a
constant applied to a $\lambda$-abstraction goes back at least to Church's simple theory of types (1940), and
λProlog's \lp{pi} and \lp{sigma} are named after Church's $\Pi$ and $\Sigma$. Computing with such terms
starts with Huet and Lang (1978) and Miller and Nadathur (1987). The phrase \enquote{higher-order abstract
syntax} comes from Pfenning and Elliott \cite{PfenningElliott1988}, and the name \enquote{$\lambda$-tree
syntax} from Miller and Palamidessi (1999). The copy clauses of \cref{sec:lambdatree-copy} go back to Miller
(1991) at second order and to Felty for arbitrary orders. The $\Llam$ fragment of the next section is
Miller's \cite{Miller1991}.

For a miniKanren programmer, the options can be compared directly. The first column is the
\enquote{functional} reading of HOAS that we ran in \cref{sec:lambdatree-beta}. The second is nominal logic
programming as in $\alpha$Kanren \cite{ByrdFriedman2007,Byrd2009}, based on nominal unification
\cite{UrbanPittsGabbay2004}. The third is $\lambda$-tree syntax.
\begin{center}
\small
\begin{tabular}{@{}>{\raggedright\arraybackslash}p{0.27\linewidth}>{\raggedright\arraybackslash}p{0.2\linewidth}>{\raggedright\arraybackslash}p{0.2\linewidth}>{\raggedright\arraybackslash}p{0.23\linewidth}@{}}
\toprule
 & host closures & names + freshness ($\alpha$Kanren) & $\lambda$-tree syntax \\
\midrule
$\alpha$-equivalence & no (\scm{eq?} only) & yes, via swapping & yes, built in \\
substitution & host application & a user relation with freshness side conditions & meta-level $\beta$ (or $\beta_0$ + copy) \\
analysing a body & only by applying to a name & by unification under \texttt{tie} & by mobility (\texttt{pi}, \texttt{=>}) \\
building an abstraction by unification & impossible & yes & yes (pattern unification) \\
\enquote{$x$ not free in $M$} & not expressible & freshness constraint $a\,\#\,M$ & scoping of variables \\
\bottomrule
\end{tabular}
\end{center}
The same notes (\book{209}--\book{210}) name nominal logic as the alternative approach: it is
first-order, with name-swapping and freshness, and binding is derived from these primitives. Porting the book's
programs to $\alpha$Kanren would mean rewriting them, not transliterating them: there is no $\beta$, mobility is replaced by fresh names, and hypothetical
contexts must be threaded explicitly.

\section{The \texorpdfstring{$L_\lambda$}{Lλ} subset}\label{sec:lambdatree-llambda}

Predicate quantification is useful for manipulating $\lambda$-tree syntax, but \enquote{it is certainly not
central to it} \book{205}. Can $\beta$ be weakened further? The book recalls the chapter's two principles: bound
names are a fiction ($\alpha$), and binders are mobile, from term level to formula level to proof level. It then
identifies the variables that carry mobility.

\begin{definition}[Essentially universal and essentially existential occurrences, \book{205}]\label{def:lambdatree-essential}
In a goal formula $G$, a bound variable occurrence is \emph{essentially universal} if it is bound by a positive
$\forall$, a negative $\exists$, or a term-level $\lambda$. It is \emph{essentially existential} otherwise, that is,
if it is bound by a negative $\forall$ or a positive $\exists$. In a program clause the classification is dual: an
occurrence is essentially universal if it is bound by a negative $\forall$, a positive $\exists$, or a term-level
$\lambda$, and essentially existential if it is bound by a positive $\forall$ or a negative $\exists$.
\end{definition}
During proof search, essentially existential variables become logic variables, instantiated by unification
with general terms. Essentially universal ones are instantiated only with eigenvariables. \enquote{It is the
essentially universal bound variable occurrences that provide mobility during proof search} \book{205}. In the
clause \lp{typeof (abs M) (arr A B) :- pi x\ typeof x A => typeof (M x) B.}, the clause variables \lp{M},
\lp{A}, \lp{B} are essentially existential, while the \lp{pi}-bound \lp{x} sits in the negative body of the clause
and is essentially universal.

\begin{definition}[$\Llam$, \book{205}--\book{206}]\label{def:lambdatree-llambda}
$\Llam$ restricts hohh by requiring (1) that all quantification be over non-predicate types, and (2) that
\enquote{Every subterm in B with an essentially existential variable occurrence as its head, that is, a subterm of
the form ($x\,t_1\ldots t_n$) ($n\ge0$), where x is essentially existentially quantified in B, must be such that
$t_1,\ldots,t_n$ are distinct essentially universal variable occurrences bound in the scope of x.}
\end{definition}
Such an application $x\,y_1\ldots y_n$ is a \emph{higher-order pattern}. If $x$ is ever instantiated with a
term $t$, the only redexes created have the form $(t\,y_1\ldots y_n)$ with the $y_i$ not free in $t$. Up to
$\alpha$ and $\eta$, $t=\lambda y_1\ldots\lambda y_n\,t'$, and the reduction just strips the binders
\book{206}. Only this sub-case of $\beta$ is needed.

\begin{definition}[$\beta_0$, \book{206}]\label{def:lambdatree-beta0}
\[
(\abs{x}{s})\,x = s \qquad(\beta_0)
\]
In the presence of $\alpha$-conversion this is equivalent to $(\abs{y}{s})\,x=\subst{s}{x}{y}$, provided $x$ is not
free in $\abs{y}{s}$.
\end{definition}

\begin{example}[Goal versus clause, \book{206}]
With $p:\typ{i}\to\typ{o}$ and $f:\typ{i}\to\typ{i}$, the formula
$\forall_{\typ{i}\to\typ{i}}x\,\forall_{\typ{i}}y\,(p\,(x\,y)\imp p\,(f\,y))$ is an $\Llam$ \emph{goal}: both
quantifiers are positive, so $x$ and $y$ are essentially universal, and nothing existential is applied. As a
\emph{program clause} it is not in $\Llam$: there both $x$ and $y$ are essentially existential, and $(x\,y)$ applies
an existential variable to an existential one. Membership in $\Llam$ depends on polarity.
\end{example}

The book states which programs of the chapter belong to $\Llam$ \book{206}: \enquote{all the program clauses in
Figures 7.7 through 7.12 belong to $L_\lambda$}, and exactly one clause of Figure~7.4 does not, namely
\lp{redex (app (abs R) N) (R N).}, where \lp{R} is applied to the existential \lp{N}. By the same definition, more
programs of the chapter are outside $\Llam$, although the book singles out only \lp{redex}. The table summarises the
situation; the classification of the rows other than \lp{redex} and Figures~7.7--7.12 is ours.
\begin{center}
\small
\begin{tabular}{@{}p{0.47\linewidth}p{0.45\linewidth}@{}}
\toprule
in $\Llam$ (applications are \texttt{(M x)} with \texttt{x} \texttt{pi}-bound) & outside $\Llam$ \\
\midrule
\texttt{term}, \texttt{bnorm}, \texttt{red1}, \texttt{reduce}, \texttt{path}, \texttt{bpath}, \texttt{addbeta},
\texttt{typeof}, \texttt{trans}, \texttt{vacuous}, \texttt{nnf}, \texttt{prenex}, \texttt{fohcG} \ldots,
\texttt{copy}, every \texttt{subst} of Figures 7.10--7.12
& \texttt{redex} (\texttt{(R N)}); \texttt{list\_instan} (\texttt{(B T)}); \texttt{interp}/\texttt{backchain}
(\texttt{(G X)}, \texttt{(D1 X)}); \texttt{cbn}/\texttt{cbv} (\texttt{(R N)}, \texttt{(R U)}); the one-line
\texttt{subst R N (R N)}; \texttt{repeat} and \texttt{foreach} (they quantify over predicates) \\
\bottomrule
\end{tabular}
\end{center}
In every non-$\Llam$ clause of the chapter, the abstraction is applied to an argument of \emph{base} type
(\lp{term} or \lp{tm}). Instantiating \lp{R : tm -> tm} with \lp{N : tm} therefore never creates a new redex. It is
plain \enquote{instantiate the bound variable}, provided \lp{R} is known when the application is formed.

A non-$\Llam$ program can often be rewritten into $\Llam$ \book{206}. With the copy clauses of Figure~7.10, the
offending clause becomes
\begin{lprolog}
redex (app (abs R) N) S :- subst R N S.
\end{lprolog}
\enquote{If we do this, we obtain a specification that is entirely within $L_\lambda$.} The translation has two
costs \book{207}. The signature must be reflected into copy clauses with $[\![\cdot]\!]^{\pm}$. And copy behaves
operationally unlike equality (the \lp{copyi} pitfall above).

Why care? Compare the two rules \book{207}:
\[
(\abs{x}{M})\,N=\subst{M}{N}{x}\quad(\beta)\qquad\qquad(\abs{x}{M})\,x=M\quad(\beta_0).
\]
A $\beta_0$ step removes one abstraction--application pair, and the result is strictly smaller and has no new
redexes. A $\beta$ step may shrink the term (if $x$ does not occur in $M$), may blow it up (if $x$ occurs many times
and $N$ is large), and may create new redexes (if $x$ occurs applied in $M$). The difference carries over to
unification, the subject of \cref{ch:hounif} \book{208}:
\begin{itemize}[nosep]
  \item unification modulo $\beta\eta$ is undecidable and may have several most general unifiers;
  \item $\Llam$ unification is \enquote{much more like first-order unification}: it is decidable, and when a unifier
        exists there is a single most general one \cite{Miller1991,Nipkow1993};
  \item \enquote{an overwhelming majority of clauses in \textquote{typical} $\lambda$Prolog programs are within the
        $L_\lambda$ fragment}.
\end{itemize}

\begin{pitfall}[$\beta_0$ is not the same as $\Llam$]
All $\Llam$ unification problems are problems modulo $\alpha$, $\beta_0$, $\eta$, but not conversely \book{208}.
Finding $F$ with $\forall x\,(F\,x\,x=x)$ needs only $\beta_0$ to justify the solutions $F:=\lambda u\lambda v\,u$
and $F:=\lambda u\lambda v\,v$. It is not an $\Llam$ problem, because $F$ is applied to the same universal variable
twice, and it has two incomparable most general unifiers. A unifier for a miniKanren must decide what to do with
such non-pattern equations: reject them, or postpone them until they become patterns.
\end{pitfall}

Finally, raising (\cref{ch:lambda}), the transformation of $\forall y\,\exists x\,G$ into
$\exists h\,\forall y\,\subst{G}{(h\,y)}{x}$ that implementations use to simplify quantifier alternation,
\enquote{preserves the class of $L_\lambda$ formulas} \book{208}: $(h\,y)$ is a pattern.

\begin{kanren}[A requirements checklist for \micro]
Gathering the boxes of this chapter, a miniKanren that runs the book's programs needs the following. Each item is
justified by a program above. The design that provides them is in \cref{ch:design,ch:module}.
\begin{center}
\small
\begin{tabular}{@{}p{0.34\linewidth}p{0.58\linewidth}@{}}
\toprule
requirement & needed by \\
\midrule
abstractions as data, equal up to $\alpha$ (and $\eta$) & every program; \texttt{(all p)}; answers like \texttt{abs W1\textbackslash\ W1} \\
application of a known abstraction to a term & \texttt{list\_instan}, \texttt{interp}, \texttt{cbn}, \texttt{cbv}, \texttt{redex} \\
eigenvariables with the scoping constraint & every \texttt{pi}; correctness of \texttt{vacuous}, $\eta$-\texttt{redex}, \texttt{prenex} \\
scoped assumed facts, possibly with logic variables & \texttt{term}, \texttt{bnorm}, \texttt{typeof}, \texttt{trans}, \texttt{copy} \\
scoped assumed rules with bodies and own $\forall$ variables & \texttt{bpath}, \texttt{addbeta}, the \texttt{trans} variant \\
pattern unification for output abstractions \texttt{(D x)} & \texttt{red1}, \texttt{path} synthesis, \texttt{nnf}, \texttt{prenex}, \texttt{copy}, \texttt{subst} backwards \\
occurs check, also through abstractions & \texttt{typeof (abs x\textbackslash\ app x x)} \\
printing binders with invented names & every answer containing an abstraction \\
\bottomrule
\end{tabular}
\end{center}
Some of this exists in the miniKanren literature. The miniKanren with symbolic constraints has an
\enquote{eigen} operator with a scope check \cite{HemannByrdFriedman2013}. $\lambda$Kanren adds \texttt{tie},
application, \texttt{all} and assumed relations to $\mu$Kanren \cite{MaLuFriedman2020}, and its successor
discusses the $\alpha$-equivalence surprises of combining higher-order unification with miniKanren
\cite{MaFriedman2021} (see \cref{ch:lkanren}). Today's \micro{} has none of these features. It does have
first-order unification with occurs check and records, which can represent de Bruijn terms, and an interleaving
search, which already changes some of the book's loops into productive streams.
\end{kanren}

\section{Summary}\label{sec:lambdatree-summary}
\begin{itemize}
  \item An object-level binder is a constant applied to a meta-level abstraction: \lp{all (x\ B)},
        \lp{abs (x\ t)}. Equality of the meta-language includes $\alpha$ (bound names are a fiction), $\beta$ and
        $\eta$ (so \lp{(all p)} is a quantifier too).
  \item Object-level substitution is meta-level $\beta$: \lp{(B T)}. Interpreters (\lp{interp}) and evaluators
        (\lp{cbn}, \lp{cbv}) need no renaming code at all.
  \item An abstraction cannot be decomposed into a name and a body. Recursion under a binder uses
        \emph{mobility}: term-level $\lambda$, then formula-level \lp{pi}, then proof-level eigenvariable, with
        \lp{=>} assuming the clauses that describe the new eigenvariable (\cref{def:lambdatree-idiom}).
  \item The eigenvariable's scoping constraint is part of the meaning of \lp{pi}. It silently implements every
        \enquote{$x$ not free in $M$} side condition (\lp{vacuous}, $\eta$, prenex). A fresh constant without it is
        unsound.
  \item Assumptions are scoped, may contain logic variables (\lp{typeof x A}), and may be rules with bodies and
        their own quantifiers (\lp{bpath}, \lp{addbeta}, \lp{trans}). Some predicates have no static clause for
        variables at all.
  \item Metatheory comes cheap: the substitution lemma for \lp{bnorm}, subject reduction for \lp{typeof}, and the
        correctness of \lp{subst} all follow from instantiating a \lp{pi} and discharging a \lp{=>}.
  \item Programs that \emph{build} abstractions (\lp{red1}, \lp{nnf}, \lp{prenex}, \lp{copy}, path synthesis,
        \lp{trans} backwards) rely on unification solving \lp{(D x) = t} by abstracting the eigenvariable out:
        pattern unification.
  \item Copy clauses, generated from the signature by type lowering $[\![\cdot]\!]^{\pm}$, specify substitution
        using only $\beta_0$. They are logically equality but operationally term generators.
  \item The relations are not functions: \lp{prenex} has five answers on the book's example, \lp{subst} run
        backwards enumerates the four existential generalisations.
  \item $\lambda$-tree syntax is HOAS with $\lambda$-terms modulo $\alpha$, a weak $\beta$ and possibly $\eta$, which can
        be analysed intensionally, unlike host-language closures. $\Llam$ restricts existential variables to
        patterns, so that only $\beta_0$ is needed and unification is decidable and unitary (\cref{ch:hounif}).
  \item On today's \micro, abstractions as closures work only as inputs, and explicit contexts with de Bruijn
        indices replace \lp{pi}/\lp{=>} at the cost of threading the context and writing capture conditions by
        hand. The requirements checklist above is the starting point of the proposal (\cref{ch:design}).
\end{itemize}
